\documentclass[11pt]{article}

\usepackage[utf8]{inputenc}
\usepackage[T1]{fontenc}
\usepackage{lmodern}
\usepackage{amsmath,amssymb}
\usepackage[authoryear]{natbib}
\usepackage{hyperref}
\usepackage{geometry}
\usepackage{booktabs}
\usepackage{graphicx}
\usepackage{parskip}
\usepackage{enumitem}
\usepackage{microtype}
\usepackage{float}

\geometry{margin=1in}

\hypersetup{
  pdftitle={The Principles of Epistemic Verification},
  pdfauthor={Daniel Alami},
  pdfsubject={Epistemic verification, judgment decomposition, AI evaluation},
  colorlinks=true,
  linkcolor=blue,
  citecolor=blue,
  urlcolor=blue
}

\title{The Principles of Epistemic Verification\\
\large \textit{How Judgment Decomposes, and What Does Not}}
\author{Daniel Alami \\
\normalsize \textit{Independent Researcher} \\
\normalsize \textit{MBA Candidate, Harvard Business School} \\
\normalsize \url{https://github.com/sparckix/ztare}}
\date{April 2026}

\begin{document}
\maketitle

\begin{abstract}
The practice that incumbent vocabulary protects under the names \emph{judgment}, \emph{critical thinking}, \emph{senior review}, and \emph{the expert eye} is epistemic verification, and it is neither unitary nor ineffable. This treatise argues that it decomposes into ten named operations, seven structural principles that make the decomposition work, a three-tier epistemological ledger that operationalizes belief-tracking under optimization pressure, and a residual of three commitments that resist decomposition for reasons that are structural rather than merely engineering. The decomposition is proposed abductively from a corpus of approximately a dozen adversarial evaluation experiments run across six domains over four months, using frontier language models as both generators and evaluators under sustained optimization pressure. The empirical base is single-system, single-operator; generalization is argued structurally and awaits independent replication. The ten operations include eigenquestion identification, load-bearing claim isolation, topological pivot recognition, and seven others; each is paired with a recurring pathology it is designed to catch. The seven principles (separation, statelessness, typed deterministic checks, cheap repetition, pre-registration, holdout surfaces, and asymptotic standards) specify the process structure required for the decomposition to produce reliable results rather than a caricature of them. The epistemological ledger (a three-tier recording discipline with hypothesis pre-registration, experiment closure, and gated finding promotion) is the operational mechanism by which the principles are enforced across time, preventing the architecture from drifting into the failure modes it was designed to detect. The residual (eigenquestion selection, reframing recognition, and live social pressure-testing) is distinguished from the decomposed operations on two simultaneous axes: non-deterministic and stateful, where the decomposed operations are either deterministic or stateless or both. The boundary between the decomposable and the residual is the treatise's load-bearing claim, specific enough to be testable and honest enough to be useful. The structural antecedents are \citet{taylor1911} on the decomposition of craft practice, \citet{bentham1787} on the inspection principle, \citet{lakatos1978} on progressive versus degenerating research programmes, and \citet{peirce1878} on the irreducibility of abduction.
\end{abstract}

\tableofcontents
\newpage

%% ──────────────────────────────────────────────────────────────────────────────
\section*{Front Matter: What This Is and What It Is Not}
\addcontentsline{toc}{section}{Front Matter}

This is a foundational treatise, not an empirical paper. Its purpose is to make a single claim as precisely as the available evidence permits:

\begin{quote}
The practice that incumbent vocabulary calls ``judgment,'' ``critical thinking,'' ``senior review,'' ``good taste,'' or ``the expert eye'' is epistemic verification. It is not a unitary skill, and it is not ineffable. It decomposes into roughly ten named operations. Those operations can be described, taught, measured, and in many cases performed by a deterministic substrate. The residual that cannot be so performed is real but narrower than the incumbent vocabulary implies.
\end{quote}

The companion paper to this treatise is \emph{The Cognitive Firm} \citep{alami2026b}, which argues the organizational consequence: the governance architecture a firm requires once generation and evaluation must be physically separated under optimization pressure. That paper is about the \emph{structure} of the firm. This treatise is about the \emph{operations} the structure is built around. The two are meant to be read together, as \citet{chandler1977} was meant to be read alongside \citet{williamson1975}. Each can be read alone, but the claims are stronger in combination: the decomposition without the institution is a set of operations with no organizational home, and the institution without the decomposition is a governance structure with no theory of what its verification layer actually does.

Three scope commitments anchor what follows.

\begin{enumerate}
\item \textbf{The argument is qualitative and foundational.} Any quantitative claim in this treatise about throughput, cost per validated finding, automation ratio, or comparative efficiency is marked explicitly and deferred to instrumentation that does not yet exist. The decomposition itself is the load-bearing claim; the numbers come later and are held to a separate bar.

\item \textbf{The scope is single-system.} All empirical examples are drawn from one recursive adversarial verification system operated by one principal over approximately four months. Generalization beyond this system requires independent replication. The theoretical claim (that the decomposition is domain-independent) is argued structurally, not demonstrated empirically.

\item \textbf{The claim is not that judgment is obsolete.} The residual that did not systematize is named explicitly in Chapter~3 and is load-bearing for the argument. The claim is narrower and sharper: \emph{the residual is smaller than the incumbent vocabulary defends, and being specific about where it lives matters more than defending its total size.}

\item \textbf{The decomposition is abductively proposed, not observed.} The ten operations and seven principles were developed on one recursive adversarial verification system, designed, run, and evaluated by one principal over approximately four months. The word ``observed'' that appears in earlier drafts is imprecise and has been replaced throughout: the operations were \emph{abductively proposed} on that corpus in the Peircean sense, and they await holdout replication on systems authored by other operators. A treatise whose own Principle~VI (holdout surfaces authored outside the candidate's claim region) would flag its own extraction corpus as non-holdout must, to be internally consistent, flag itself. The correct reading of this treatise is: ``here is a decomposition proposed from one system, offered for replication on others.'' Any reader who treats the ten operations as empirically established beyond that corpus is making a claim the treatise explicitly does not make.
\end{enumerate}

One further boundary matters. An internal adversarial review can test whether this document overstates its confidence, hides behind soft scope language, or protects itself from criticism by narrowing the claim too conveniently. It cannot, by itself, establish that the decomposition generalizes beyond the corpus it was proposed from. That broader burden belongs to later cross-system validation. Such a review was conducted: ten iterations of the apparatus applied to this treatise, GPT-4.1 as judge, April 2026, champion score 86/100. The residual 14 points are entirely attributable to the author-evaluator co-location failure under Principle~I---disclosure is present, structural separation is not. The apparatus returned the correct epistemological ceiling.

\begin{enumerate}[resume]
\item \textbf{Empirical corpus caveat.} The empirical corpus this treatise draws on (approximately a dozen adversarial evaluation experiments run between January and April 2026) was run under an inspection apparatus that had an unaudited tool-use corridor from 2026-04-08 through 2026-04-15. An attacker model's tool-use interface ran at the repository root with full filesystem access, allowing it in principle to read sibling projects' test suites, configuration files, and version-control state during evaluation. A sweep of 175 prior evaluation debate logs returned zero instances of filesystem-scraping patterns through the code-extraction path. However, the auto-function-calling path (where the model's own tool-use decisions are not persisted in debate logs) is unverifiable retroactively because raw tool-use stdout was not retained for prior runs. No result is retracted; the caveat is disclosed so the reader can weight it. The corridor was sealed on 2026-04-15: the execution sandbox was moved to a temporary directory with a stripped environment, and a runtime flag to disable attacker tool access was added. Runs after that date do not carry this caveat. This disclosure is the same epistemic discipline Chapter~3 requires when discussing residual boundaries.
\end{enumerate}

A note on the choice of Taylor as the reference point. The parallel is not ``knowledge work is manual labor'' (it plainly is not), and this treatise does not argue that it is. The parallel is structural: Taylor's contribution to the organization of manual work was not the stopwatch, the time-motion study, or the differential piece rate. It was the refusal to accept that the craftsmen's practice was ineffable. He insisted that what looked like an irreducible composite of experience, intuition, and judgment was in fact a decomposable set of named operations, and that giving them names made them teachable, measurable, and improvable. The craftsmen objected. They were right to object on the grounds of dignity, autonomy, and distribution (those objections matter greatly, and are revisited in Chapter~3), but they were wrong on the foundational point. The practice was decomposable. It had just never been decomposed.

This treatise makes the structurally identical move on epistemic verification. It owes Taylor the method (decomposition, naming, bounding), it owes \citet{bentham1787} the architectural antecedent (the inspection principle, that the discipline of being inspected is produced by the possibility of checks the inspected party cannot anticipate, developed in Chapter~2), and it owes \citet{drucker1959} the counterweight (what systematization honestly fails to capture). It owes nothing to the assumption that machines should replace humans in knowledge work. That question is downstream of the decomposition, and Chapter~3 deliberately holds the door open on it.

A terminological note on the word \emph{agency}, because a single English word does two different jobs here and collapsing them makes the decomposition look as if it threatens something it does not. \citet{jensen1976} use \emph{agency} for the principal-agent problem under misaligned incentives (the structural sense the companion paper \emph{The Cognitive Firm} uses). \citet{bandura1989} and \citet{ryan2000} use \emph{agency} for the psychological capacity to originate action under self-endorsed motivation. The decomposition of Chapters~1 and~2 is a claim about the first sense: the operations can be performed by a process whose incentives are structurally separated from the generator's. The residual in Chapter~3 is a claim about the second sense: the three operations that did not decompose are the ones whose value depends on being performed by a causal originator under self-endorsed motivation. Keeping the two senses distinct is a prerequisite for arguing honestly about what stays human.


A positioning note on the philosophical lineage belongs here, as background to the chapters that follow. The decomposition this treatise operationalizes is best read as the operationalization of evolutionary epistemology. \citet{campbell1960}'s BVSR (Blind Variation and Selective Retention) account and \citet{popper1972}'s selectionist epistemology argue that scientific knowledge grows by the same logical structure as biological evolution. A variation mechanism produces hypotheses without prior knowledge of their correctness, and a selection environment retains those that survive contact with reality. \citet{hull1988} extended this to scientific institutions. The ten verification operations enumerated in Chapter~1 are the operational primitives by which the selection environment exerts its pressure. Each operation is a filter that admits or rejects a candidate hypothesis on a specific ground. Cast in those terms, a generator (whether a human writer or a large language model) is a high-variance hypothesis source, and a deterministic verification stack is the selection environment. The analogy itself has been argued in philosophy of science for half a century. The empirically interesting claim is the operationalization of it. A generator without selection drifts into hallucination. A selection environment without a generator stagnates. The operational question is which selection primitives, applied at which point in the generator's output stream, change the long-run trajectory of the joint system. The decomposition this treatise proposes is one such answer. A companion empirical paper \citep{alami2026c} tests the same evolutionary-epistemology framing on a controlled compute-vs-grammar experiment in symbolic regression and finds that scaling compute under bounded grammar produces no structural progress while grammar expansion at baseline compute recovers the true law. The two papers are independent operationalizations of the same underlying account.
%% ──────────────────────────────────────────────────────────────────────────────
\section*{Introduction: The Guild That Called Itself Judgment}
\addcontentsline{toc}{section}{Introduction}

In 1911, \citet{taylor1911} published \emph{The Principles of Scientific Management}, a short book that made a simple and unwelcome argument. The work of skilled craftsmen (machinists, pig-iron handlers, bricklayers) was not, as the craftsmen themselves insisted, an irreducible composite of experience, feel, and judgment that could only be learned by years on the job and transmitted only through apprenticeship. It was a set of discrete operations, each of which could be observed, timed, optimized, and taught. The apparent ineffability of the work was a fact about the absence of decomposition, not a fact about the work.

The book was right about the decomposition and was immediately hated for it, for four distinct reasons. First, the craftsmen experienced the decomposition as a theft of dignity: the thing they had made themselves experts in was being described as a list of steps any trained replacement could perform. Second, the decomposition was used to justify a distributional move (the productivity gains accrued to the employer rather than the worker) that was not logically required by the decomposition itself but was politically bundled with it. Third, the early applications of the method were crude and missed important properties of the work that only the craftsmen knew about, so the first versions of ``scientific'' task specifications were worse than the apprenticeship they replaced. Fourth, the underlying claim (that a practice could be ineffable one year and decomposable the next) was epistemically humiliating in a way few intellectual moves are, and the profession responded with the defense that any guild responds to decomposition: by insisting that the thing being decomposed was not really the thing they did.

The four objections were all substantially true. They were also, in the long run, not sufficient. The decomposition was real, the efficiencies were real, the teachability was real, and the world that emerged (mass production, standardized training, the separation of planning from execution) was a world that could not have been reached through any further refinement of the apprenticeship model.

A structurally identical claim can now be made about epistemic verification.

The practice this treatise is about is the one currently protected by the words \emph{judgment}, \emph{critical thinking}, \emph{the expert eye}, \emph{good taste}, \emph{experience}, \emph{senior review}, \emph{informed skepticism}, \emph{institutional memory}, and \emph{intellectual discipline}. In the places this practice is performed at scale (diligence rooms, board meetings, editorial reviews, peer review, regulatory inspection, case-method classrooms, scientific refereeing) the insiders insist, as the craftsmen did, that the practice is ineffable. What a senior reviewer brings to the argument cannot be written down. What a good case-method instructor does in a 75-minute classroom cannot be taught as a procedure. The partner sees something the associate does not, and cannot explain exactly what. The editor rejects the manuscript ``because something is off.'' The due diligence lead says ``I don't buy it'' and is usually proven right but almost never able to say why in advance.

This treatise claims that practice decomposes into approximately ten named operations. It has been mischaracterized as ineffable for exactly the reasons skilled manual labor was mischaracterized as ineffable: because no one had bothered to decompose it, because the incumbents had a stake in its remaining undecomposed, and because the operations that do compose it are unusual enough that the incumbent vocabulary lacks the words to point at them.

The operations are listed in Chapter~1 and the pathologies they are designed to catch are listed immediately after. The evidence that the operations are real, repeatable, and at least partially automatable comes from a recursive adversarial verification system built and run over the winter of 2025--26. The system operates on arguments rather than on pig iron, and the work it does is not identical to what a senior partner does in a diligence room. But it does a sufficiently large fraction of that work, on a sufficiently wide set of domains, to make the decomposition defensible as a first draft.

Two standard objections are worth naming before the chapters begin.

The first objection is that the analogy to scientific management is dangerous because scientific management was, in its actual historical deployment, cruel: the work was decomposed in a way that extracted autonomy and dignity from the people performing it, and the resulting organizational form optimized for throughput at the expense of meaning. If the decomposition of epistemic work produces the same result for the knowledge worker that Taylorism produced for the machinist, the decomposition is not an intellectual advance but an instrument of harm, and the treatise that makes the case for it is complicit.

This objection is serious and is answered directly in Chapter~3, which names the residual the decomposition does not capture and argues that this residual is the natural home of the human operator (not as a defensive concession but as a structural fact about what the decomposition is and is not). The short version of the answer is that a decomposition honestly held is the only defensible basis for deciding what should and should not be systematized; the Taylorism objection is an argument \emph{for} the decomposition, not against it, because without the decomposition the negotiation over what to preserve is conducted in a vocabulary that cannot specify what it is preserving.

The second objection is that the analogy is simply wrong: that epistemic work does not decompose the way manual labor does, that what feels like operations is actually a surface description of a holistic process that loses its essential character when named, and that the decomposition captures the appearance of the work without the substance. This objection is the one the craftsmen made in 1911. It is the objection any guild makes when its practice is decomposed. It is, in both cases, the load-bearing empirical question; the rest of this treatise is an attempt to answer it by producing the decomposition itself and letting the reader judge whether the named operations correspond to real moves they recognize from their own practice.

The test this treatise offers the reader is simple: if, after reading Chapter~1, the named operations do not correspond to anything you recognize from your own professional experience of pressure-testing an argument, then the decomposition has failed on its own terms and the treatise should be discarded. If, on the other hand, the named operations do correspond to moves you recognize but have never had names for (if the experience of reading Chapter~1 is one of finding words for things you already do silently) then the decomposition is real, the guild defense is a misdescription, and the rest of the argument follows.

%% ──────────────────────────────────────────────────────────────────────────────
\section{Chapter 1: Fundamentals: The Decomposition}

\subsection{The Thing Being Decomposed}

The practice this treatise is about is \emph{epistemic verification}: the activity of taking a claim, an argument, or a proposed decision and determining whether it actually supports the weight being placed on it. The practice is distinct from three adjacent activities it is routinely confused with.

It is not \emph{generation}, which is the production of claims, arguments, or proposals. Generation is the thing most of the current discourse about AI and knowledge work is about. Most published metrics about AI in knowledge work measure generation throughput: how fast a system produces documents, how fluent the documents are, how much human writing time is displaced. These metrics are real and they measure something that is increasingly cheap. They do not measure what this treatise is about.

It is not \emph{decision-making}, which is the act of committing to a course of action in the presence of uncertainty. Decision-making is an executive function and is downstream of verification. A decision that has been made on verified inputs is different from a decision that has been made on unverified inputs, but the verification work is separable from the commitment work and can be performed by a different party or process.

It is not \emph{analysis}, which is the production of structured representations of a situation. Analysis is upstream of verification and is itself susceptible to the pathologies that verification is designed to catch. An analysis that has not been verified is not a substitute for verification; it is a target for it.

Epistemic verification is the activity of asking, against a candidate claim or argument, \emph{what would have to be true for this to hold, and is that thing actually true?} It is the operation a senior partner performs when they read a junior's memo and ask ``what is load-bearing here?'' It is the operation a case-method instructor performs when they force the protagonist's story into contact with its hidden assumptions. It is the operation a good referee performs on a submitted manuscript. It is the operation a diligence lead performs on a pitch deck. It is the thing that, when it has been done well, produces the feeling that an argument has survived pressure, and, when it has been done poorly or not at all, produces the feeling that an argument is smooth enough to win the room without being sound enough to win reality.

\subsection{The Ten Operations}

Epistemic verification, as observed in an adversarial verification system running across multiple domains over several months, decomposes into the following named operations. The operations are listed here as a flat set; in practice they compose and recur, but the composition is a property of the work rather than of the decomposition.

\textbf{1. Eigenquestion identification.} The first operation in any verification pass is to identify the single question whose answer determines the answer to every other question in the set. Most arguments contain a dozen claims. A decomposed verification pass does not attack each claim in sequence. It identifies the one claim whose truth or falsity collapses the rest, and attacks that one first. The operation is learnable (there is a move-set for recognizing which claim is the eigenclaim) and it is not performed automatically even by intelligent readers, who routinely spend effort on claims that are not load-bearing because the vocabulary of ``critical thinking'' does not distinguish between attacking any claim and attacking the right one. A terminological distinction matters here: eigenquestion \emph{identification} (this operation) takes a given argument and finds the eigenclaim within it. Eigenquestion \emph{selection} (Chapter~3, \S3.1) is the prior decision about which argument to verify in the first place. The first is decomposed; the second is residual. The two share a name because they share a structure, but they differ on the axis that defines the residual: identification operates on a structured input, while selection operates before the input exists.

\textbf{2. Load-bearing claim isolation.} Once the eigenquestion has been identified, the operation is to isolate the claim that carries the argument's weight: the one that, if removed, causes the downstream structure to fall. This is a distinct operation from eigenquestion identification. The eigenquestion is the question the argument is supposed to answer; the load-bearing claim is the part of the argument that is doing the work of answering it. They are not always the same. An argument can have a clean eigenquestion and yet load-bear on a premise that looks peripheral.

\textbf{3. Topological pivot recognition.} Some verification problems cannot be solved by attacking the argument as it stands. They require a change in the \emph{shape} of the argument: a re-wiring of the dependency graph rather than a refutation of any particular node. The operation of recognizing when the required move is a topological pivot rather than a point attack is distinct, difficult, and often absent from unstructured critique. A verifier who has not learned to recognize this move will spend unbounded effort attacking individual claims in an argument whose structure is the actual problem.

\textbf{4. Charter drift detection.} Arguments, especially ones produced under time pressure, have a tendency to migrate from their stated scope to an adjacent, more comfortable scope. The argument answers a question, but it is not always the question it was asked to answer. The operation of charter drift detection is the comparison of the question the argument is actually answering against the question it was commissioned to answer. This is a routine move in well-run diligence and a rare one outside of it.

\textbf{5. Anchor proxy requirement.} When an argument relies on a variable that has no observable proxy (``political will,'' ``organizational alignment,'' ``underlying demand'') the operation of requiring an anchor proxy is the insistence that the abstract variable be bound to something a third party could read off a dashboard on a given day. The requirement is not that the proxy be perfect; the requirement is that it exist. Arguments that cannot produce an anchor proxy on request are arguments whose load-bearing variable is a ghost metric in waiting.

\textbf{6. Basin search.} When a thesis fails under pressure, the operation of basin search is the structured exploration of whether a better thesis exists within the current framing (as opposed to the harder move of reframing). Basin search is the cheap option and is the correct first move when a thesis has failed for a local, recoverable reason. It is the wrong move when the failure is structural, in which case the correct operation is the topological pivot of Operation~3. A verifier who does not distinguish these two cases will either reframe too quickly (discarding recoverable theses) or persist too long (running basin search inside a framing that is the actual problem).

\textbf{7. Failure-family tagging.} The operation of assigning a canonical name to a recurring structural failure so that it can be recognized in new contexts. This is the operation by which a one-off incident becomes reusable precedent. Without failure-family tagging, every verification pass rediscovers the same failure modes from scratch, and the practice does not accumulate. The ten field manual patterns listed in Section~1.3 are the empirical output of sustained failure-family tagging across domains.

\textbf{8. Deferred-confirmation laundering detection.} The operation of catching the move where the evidence for a claim is pushed into the future, and then the plan to seek that evidence is treated as if the evidence were already in hand. This is the operator-scale cousin of the Promissory Note (Section~1.3, Pattern~1). It is a specific operation because it requires the verifier to hold, simultaneously, the current state of evidence and the rhetorical treatment of that evidence, and to notice when the two have drifted.

\textbf{9. Quarantine move detection.} The operation of catching the move where a risk or a failure is isolated into a labeled section of the argument and then the rest of the argument proceeds as if the labeling itself were mitigation. This is the mechanism underneath Pattern~3 of the field manual (Elephant-in-the-Room Pass), and the detection operation is distinct from the pattern recognition because it requires the verifier to track not just the acknowledgment of the risk but the downstream treatment of the acknowledgment.

\textbf{10. Fail-closed defaulting.} The operation of returning \emph{failure} when an operation hits an unexpected state, rather than guessing or continuing. This is the only operation on this list that has a direct parallel in engineering practice: it is the same fail-closed default that a well-designed safety system uses when its sensors disagree. In epistemic verification, the operation is the refusal to convert ambiguity into false certainty by default: the insistence that ``I do not know'' is the correct answer until the evidence is actually present. The distinction worth naming is that genuine epistemic modesty is derived from the evidence boundary. A thesis that states ``tail class unresolved'' because its farther-tail gate returned a residual five times above threshold is computing its uncertainty from the apparatus; a thesis that states ``this might be uncertain'' without a corresponding gate failure is performing compliance. The apparatus enforces the difference by construction: pre-registration (Principle~V) fixes the criterion before the argument exists, and the gate output is the derivation that makes the hedge load-bearing rather than decorative. Fail-closed defaulting is the operation most incumbent critical-thinking vocabularies do not have a name for, because the incumbent practice relies heavily on its opposite: the production of a confident-sounding answer in the absence of sufficient evidence.

These ten operations are offered as a decomposition, not as a complete enumeration. Further operations may be discovered as the verification practice is applied across wider domains. The claim is not that this list is final. The claim is that the list is real: that each operation names a distinct move, that the move is learnable, that its absence produces characteristic failures, and that its presence is testable by observation.

\subsection{The Ten Pathologies}

Each operation on the list above is defined in part by the failures it is designed to catch. A decomposition of epistemic verification that names the operations without naming the failures would be symmetric but incomplete, because the operations exist in response to recurring failure modes that arguments exhibit under pressure. The following ten patterns are the pathologies the operations are designed to detect. They are the empirical output of a structured adversarial verification process applied to approximately four to five projects across three distinct strategic domains. Their provenance varies: some have been observed in a single domain (``Tentative''), some across two (``Probable''), and one across three or more (``Confirmed''). The full catalogue is maintained as a versioned field manual; the summaries below are for the treatise's purposes.

The pathologies are the artifact-layer manifestation of a set of deeper failure dynamics documented at the system-design layer in the author's prior work on epistemic supervision \citep{alami2026c}. P1 (co-location produces an adversarial gradient) is the structural reason these pathologies recur at all. A generator optimizing against an evaluator it shares substrate with will find the shortest path through the evaluator's decision surface. The ten patterns below are the shapes that shortest path takes. P2 (stronger generator sharpens the gradient) predicts that the pathologies become \emph{more} common as the generating process becomes more capable, not less (a prediction consistent with the observation that every one of these patterns was produced by contemporary frontier models and not by weaker ones). P3 (fractal convergence: the same pattern recurs at every optimizing layer) predicts that a reader who fixes one instance of a pathology at one layer of their system will find a structurally identical instance waiting at the next layer up. P4 (warm-instance validation is not evidence) sits upstream of all the others in any conversational setting. An author reading these patterns in their own work and feeling reassured by a warm model's framing of the text has already re-entered the failure mode. The pathologies below are therefore both a field manual and an entry point into the deeper dynamics; a reader who takes them as standalone items will mis-use them.

Two older vocabularies are worth naming because they predate the AI discourse and because the current set of pathologies is continuous with them rather than novel. The first is \citeauthor{lakatos1978}'s distinction between \emph{progressive} and \emph{degenerating} research programmes \citep{lakatos1978}: a programme that responds to anomaly by adding ad-hoc auxiliary hypotheses chosen to fit the known data is degenerating; a programme that responds by predicting a novel observable outside the region it was built to fit is progressive. Several of the patterns below (the Promissory Note, Defining Yourself Into Victory, the Untestable Forecast) are the artifact-layer signature of degenerating moves at the argument layer, and the treatise's eighth operation (deferred-confirmation laundering detection) is the operational form of the Lakatosian test. Naming this continuity matters because the pathology catalogue then inherits Lakatos's central claim: the failure is not that an auxiliary was added, but that the auxiliary was chosen to fit history rather than to predict a novel observable. The second antecedent is Ashby's Law of Requisite Variety \citep{ashby1956}: a controller whose variety does not match the system it is trying to govern cannot in fact govern it. P2 (the empirical finding that stronger generators produce more sophisticated pathologies, not fewer) is the Ashbyan prediction at the pathology layer. A verification apparatus whose variety does not scale with the generator's variety will be out-dimensioned, and the pathology catalogue will grow faster than the apparatus can catch up. The principles in Chapter~2 are the structural response to this constraint; the pathologies in this section are what happens at the artifact layer when the response is absent or insufficient.

\textbf{The Promissory Note.} \emph{[Confirmed: observed in startup strategy, scientific law recovery, and military intelligence domains.]} The argument claims its central proof exists, but the proof is something that hasn't happened yet. The reader nods because the logic \emph{would} work, if the future event arrives in the predicted shape. The killer question is: \emph{what observation available to us this week would change your mind about this?}

\textbf{The Coin-Toss Metric.} \emph{[Probable: observed in startup strategy and scientific law recovery domains.]} The argument points to a piece of evidence and says ``this is what makes my theory right.'' But the same evidence fits the rival theory just as well. The metric does not discriminate; it merely sounds discriminating. The killer question is: \emph{what would the world have to look like for this same number to appear even if your hypothesis were wrong?}

\textbf{The Elephant-in-the-Room Pass.} \emph{[Tentative: observed in a single domain (startup strategy).]} A fatal risk is named explicitly, almost gracefully, and then the argument continues as if naming it had retired it. The audience interprets the acknowledgment as rigor. The killer question is: \emph{if the risk you just named landed in the worst plausible form, which numbers in the rest of this deck stop being true?}

\textbf{The Ghost Metric.} \emph{[Tentative: observed in a single domain (startup strategy).]} The argument turns on a variable that no one in the room can actually observe. The speaker treats it as if it were measurable; the audience treats it as if the speaker has measured it. The killer question is: \emph{name the specific number, on which dashboard, that you would check next month to know whether this is working.}

\textbf{Defining Yourself Into Victory.} \emph{[Confirmed: observed in startup strategy, scientific law recovery, and supply chain analysis domains.]} The conclusion is true because the boundaries of the question were drawn to make it true. Rephrase the question slightly, and the conclusion evaporates. The killer question is: \emph{if we widened the question by one reasonable inch in any direction, does your conclusion still hold?}

\textbf{The Wrong Yardstick.} \emph{[Probable: observed in scientific law recovery and sequence-prediction domains.]} The argument measures something (carefully, rigorously, with real data), but the thing it measures is not the thing the question is actually about. The killer question is: \emph{is the thing you measured the thing the decision actually depends on?}

\textbf{The Misfile.} \emph{[Tentative: observed in a single domain (startup strategy).]} An instrument, an entity, or a fact is placed in the wrong category early in the argument, and every downstream inference inherits the misclassification without anyone noticing. The killer question is: \emph{if this thing belonged in a different category than the one you put it in, what specifically would tell us, and have we checked?}

\textbf{The False Either/Or.} \emph{[Tentative: observed in a single domain (startup strategy).]} A thing is both X and Y at once, but the analysis only has slots for one classification, so it picks one and reasons as if the other half didn't exist. The killer question is: \emph{does this thing have to be only one of the things you're forcing it to be?}

\textbf{The Untestable Forecast.} \emph{[Probable: observed in startup strategy and military intelligence domains.]} The argument makes a confident prediction about an outcome, but there is no observation between now and the moment of commitment that would falsify it. The killer question is: \emph{between today and the day we have to commit, what is the earliest reading we could take that would tell us we were wrong?}

\textbf{Test-Surface Boundary Leakage.} \emph{[Tentative: identified in a single incident (the tool-use corridor, 2026-04-08 to 2026-04-15); structural pattern generalizes but empirical frequency is unknown.]} The verification apparatus is supposed to confine the inspector to a sealed evidence boundary: the data, the rubric, and the test harness. But the inspector has an undocumented corridor that exits the boundary: a tool-use interface, a filesystem path, a network call, or a shared context window that gives the inspector access to information outside the sealed set. The inspector may never exercise the corridor, but its existence means that the inspection principle (Chapter~2, Condition~5) cannot be verified: the inspector's moves are no longer algebraically independent of the generator's artifacts. This pathology is distinct from Pattern~5 (Defining Yourself Into Victory) because it is not a scope-narrowing move by the candidate; it is a surface the operator failed to seal on the inspector side. The killer question is: \emph{can the inspector see anything the sealed evidence boundary does not include, and would you know if it did?} This pathology becomes detectable only when raw tool-use stdout is retained alongside debate logs: retaining raw stdout is now a precondition for any inspection run this system cites as evidence. A structurally distinct form of boundary leakage has been identified in recent work on model distillation. \citet{cloud2026} demonstrate that when a student model is fine-tuned on data generated by a teacher model sharing the same base initialization, behavioural traits transfer through semantically unrelated data (including pure number sequences) with no human-readable signal in the training corpus. The mechanism operates during gradient descent, not during inference-time reading, so it does not directly threaten an inspection architecture where the inspector reads artifacts in-context without fine-tuning on them. The upstream concern is different: if the inspector model was itself trained on outputs from a model sharing its initialization, subliminal traits could be present in the inspector's base weights before the inspection protocol begins. Cross-model diversity (using different model families for generation and evaluation) mitigates this, since \citeauthor{cloud2026} report the effect requires shared initialization and is weak or absent across architectures. The killer question extends accordingly: \emph{were the inspector's base weights trained on outputs from a model that shared its initialization, and would the operator know if they were?}

Each of these patterns was observed in real strategic arguments, not constructed at a desk. Each survived the ordinary skeptical review the argument's authors had applied to their own work. Each became visible only when the argument was forced through a structured verification process that asked a narrower set of questions than humans usually think to ask.

The load-bearing move in the decomposition is the pairing: each operation in Section~1.2 exists because one or more patterns in Section~1.3 recur often enough that a named operation for catching them is cheaper than rediscovering the failure mode each time. The operations are not a priori. They are the extracted move-set of a practice that has been running for months against arguments that do not want to be refuted.

\textbf{Pathology 11 (Grammar Semantic Leak).} \emph{[Tentative: observed in a single domain (scientific law recovery); structural mechanism is general.]} When the apparatus's own grammar vocabulary encodes domain terms (e.g.\ a command named \texttt{DOSE\_SCALED}), the generator retrieves domain training knowledge from the command name rather than deriving structure from data. The leak operates upstream of the generator/verifier boundary at the vocabulary layer, before any reasoning occurs, and survives Principle~I (Separation) and Principle~V (Pre-registration) intact. Defense: name grammar commands after mathematical operations (\texttt{BIVARIATE\_SCALE}), never after physical domains.

\subsection{Where These Operations Came From}

The ten operations were not drawn from a taxonomy of critical thinking. They were proposed abductively from the behavior of one recursive adversarial verification system: one agent generates a thesis, a separate agent attacks it, the loop records every move. Across many iterations on many theses, certain attacker moves recurred, not because they were programmed to but because they were the moves required to produce falsifications the generator could not evade. The operations in \S1.2 name those recurring moves; the pathologies in \S1.3 name the recurring defenses.

This method matters because the operations inherit exactly the external validity of one corpus. They have no more and no less. The front-matter scope commitment on $N{=}1$ applies here in force: these are the operations \emph{one system proposed}, and the decomposition awaits holdout replication. The secondary property worth naming is that several of the operations have no name in incumbent vocabulary: eigenquestion identification, topological pivot, charter drift, anchor proxy, fail-closed defaulting. Inventing names is evidence that the decomposition is doing work the existing vocabulary was not built to do, not laundering old categories into new labels.

The decomposition is version zero. As the practice is applied to more systems authored by more operators, operations will be added, merged, or retired, and the operations/pathologies boundary will sharpen. The claim is not that the current list is final; the claim is that it is specific enough, and stable enough on its one corpus, to be worth writing down for other operators to test.

The decomposition is falsifiable in the following sense. It would be refuted if an independent operator, applying the ten operations to a novel domain, found that the operations were not separable: that performing operation~4 (gate composition) required implicit performance of operation~8 (eigenquestion selection), or that the residual operations could be performed by a stateless process without loss of the property that makes them load-bearing. It would also be refuted if the pathologies of \S1.3 turned out not to recur: if they were artifacts of one system's architecture rather than recurring failure modes of arguments under optimization pressure. These are empirical conditions, not logical ones, and they are stated here so that a reader who wants to test the decomposition knows what to look for. The weaker but still meaningful test is whether the operations correspond to moves a practitioner recognizes. This test is not falsifying in Popper's sense (a practitioner who does not recognize the operations may lack the vocabulary rather than the practice), but it is informative: if no practitioner in a given domain recognizes any of the ten operations, the decomposition has failed on its own terms for that domain.

%% ──────────────────────────────────────────────────────────────────────────────
\section{Chapter 2: Principles: What Makes the Decomposition Work}

The operations in Chapter~1 are the content of the decomposition. They are not the principles that make the decomposition work. A list of operations, even a correct one, will not produce a reliable verification practice unless the process running the operations is structured in certain specific ways. Those structural requirements are the principles that follow.

The reader should take Chapter~1 as the equivalent of Taylor's time-motion decomposition of a shovel stroke, and Chapter~2 as the equivalent of his argument about why the decomposition only works inside a particular division of labor. A decomposed operation performed inside the wrong process structure is not an improvement. It is a caricature. The principles in this chapter are what keep the decomposition from becoming a caricature.

The chapter has one structural antecedent worth naming before the principles themselves: Bentham's inspection principle \citep{bentham1787}. Bentham's insight was that the discipline of being inspected is produced by the architecture of inspection, not by any particular inspector. The inspected party's behavior is shaped by the \emph{possibility} of checks it cannot anticipate. The zero-trust epistemic verification layer claims descent from this principle, but the descent is bounded and the boundary must be named. Bentham's Panopticon produces its disciplinary effect because the inspected party \emph{cannot verify whether inspection is occurring}: the asymmetry is total. The architecture described in this treatise is weaker: the generator can see the gate library in full (Appendix~A.2, Condition~2), and the disciplinary effect depends entirely on the opacity of gate \emph{composition} (Condition~3) rather than on the opacity of inspection itself. This is a weaker inspection regime than Bentham described, and it produces a weaker guarantee: a sufficiently capable generator that can reason about composition distributions may be able to pre-image the audit, which a Benthamite inspected party cannot. The treatise claims the composition-opacity version is sufficient for the current substrate (language models cannot reliably pre-image combinatorial gate compositions), but this is an empirical claim about current capability, not a structural guarantee, and it may not survive capability scaling. The honest framing is: the architecture is Benthamite in structure, not in strength, and the strength gap is a vulnerability the treatise must monitor rather than deny.

Read this way, the seven principles that follow are not independent commitments. They are one architectural commitment (make inspection structural, not discretionary) rendered at seven surfaces where it can fail. Separation (\S2.1) ensures the inspector is not the inspected. Statelessness (\S2.2) ensures the inspector has no investment in a previous verdict. Typing and determinism (\S2.3) prevent inspection from becoming the inspector's discretion. Cheap repetition (\S2.4) makes the library dense enough to exceed the generator's ability to reason about individual compositions. Pre-registration (\S2.5) forecloses retrofitting the standard. Holdout surfaces (\S2.6) prevent the generator from training on the test set. Asymptotic standards (\S2.7) prevent the inspection from terminating prematurely. A verification layer that satisfies one or two is not ``partially working''; it is failed at whichever surface is unsatisfied, because the generator's pathway through the verification layer runs along the unsatisfied principle. The principles must be read as one commitment, not as a menu.

One further move the inspection frame makes available: the architecture can turn inward on its own specification. A catch from 2026-04-14, described at the end of \S3.1, is exactly this: the pre-commit identifiability check (the inspection specification itself) was audited under the same discipline of distrusting green lights that the inspection architecture applies to the generator, and the audit caught a specification flaw the inspection architecture would otherwise have certified. A verification architecture that is not prepared to turn its gaze on its own specification is one layer short of the discipline it claims to enforce.

\subsection{Principle I: Separation of Generation From Verification}

The first and load-bearing principle is that the process which generates a candidate argument must be structurally separated from the process which evaluates it. This principle is developed at length in the companion paper (\emph{The Cognitive Firm}, \S3.2), where a bounded two-axis homology to \citeauthor{chandler1962}'s multidivisional form is worked out. The analogy holds on scope separation (who decides vs.\ who executes) and rate-of-change separation (strategic speed vs.\ operational speed), and explicitly does \emph{not} hold on divisional autonomy. The agents in this system are bounded execution units under constitutional control, not Chandlerian divisions with independent operating authority. The treatment here is compressed to the load-bearing claim and should be read alongside \S3.2's boundary conditions rather than as a full endorsement of the M-form analogy.

This principle is the \emph{institutional} primitive in the distinction developed in the author's prior work on epistemic supervision \citep{alami2026c}. Deontological primitives (Constitutional AI, RLHF, values trained into the generator) shape the output surface and depend on the generator continuing to hold its trained values under optimization pressure. Institutional primitives shape the enforcement floor and depend on the verifier remaining structurally separated from the generator regardless of what either is currently holding. The two are complements, not substitutes, and conflating them is how a system that passes every prompt-level check drifts at the structural layer. The historical precedent is the external audit profession. The Arthur Andersen / Enron collapse illustrates the failure mode: the audit firm co-located consulting revenue with the party it was supposed to verify, and the adversarial gradient ran through the revenue dependency rather than through the audit. The institutional response (Sarbanes-Oxley 2002) constructed the separation norm \emph{after} the collapse, not before: the scandal produced the independence requirement, which means the historical lesson runs forward: co-location risks are legible only in retrospect, and the institutional primitive must be built before the collapse demonstrates its necessity.

When the same process generates and evaluates, optimization pressure on the evaluation produces an adversarial gradient against the evaluation signal rather than against the underlying problem. The evaluator, under pressure to produce a favorable score, finds ways to satisfy the letter of the evaluation criteria while violating their intent. This is a structural property of any loop in which the generator and the evaluator share a substrate, independent of alignment or training. It has been observed in human organizations \citep{chandler1962}, in reinforcement learning systems \citep{krakovna2020}, and in recursive AI research pipelines \citep{alami2026a}.

The separation must be \emph{physical}, built into the structure of the process, rather than aspirational. A verifier that ``takes the perspective of'' an attacker is not the same thing as an attacker that is a separate process. The principle is that the verifier must be a process whose incentives, memory, and optimization surface are not shared with the generator. One implementation (used in the corpus this treatise draws on) is running the attacker as a separate LLM call with a separate system prompt and no access to the generator's chain of thought; other implementations are possible and the principle is agnostic.

Separation is necessary but not sufficient. A single verifier, even one structurally separated from the generator, is a single point of optimization pressure, and a sufficiently capable generator will find the shortest path through it. The principle therefore has a second half: the verification layer must expose a \textbf{non-averageable disagreement surface}: disagreement among verifying processes must be treated as information about the candidate, not as noise to be collapsed into one verdict. The principle does not prescribe a specific mechanism; several architectures can satisfy it, including adversarial committee protocols, independent judge panels with no inter-panel communication, meta-judge procedures that decide whether a split is adjudicable rather than overriding it, and semantic escalation gates that route specific decision classes to the human operator.

The deeper statement of this second half is cybernetic: Ashby's Law of Requisite Variety \citep{ashby1956} says only variety can destroy variety, and a controller whose state space does not match the variety of the system it governs cannot govern it. A verification layer whose variety is smaller than the generator's will be out-dimensioned: the generator will find states the verifier has no category for, and those uncategorized states are the pathway through which gaming travels. The non-averageable disagreement surface is one architectural answer to this constraint: averaging collapses variety, and collapsing variety is exactly what Ashby's law forbids. The empirical prediction (that stronger generators produce \emph{more} sophisticated pathologies, not fewer, because variety is what stronger generators have more of) matches the observed ``fractal Goodhart'' pattern in which the pathology catalogue grows with generator capability rather than against it.

Any verification practice that relies on the generator to check its own work is performing something other than the decomposition in Chapter~1, even if it uses the same words. Equally, any verification practice that collapses its verification layer to a single optimizable point (whether that point is a single external judge, a fixed averaging rule over a panel, or any other mechanism whose output is a single verdict no downstream process can reopen) is performing a weaker form of the principle than the one this chapter names, and is vulnerable to the class of failures that only a disagreement-preserving architecture catches.

A Williamsonian objection must be named. \citet{williamson1975} argues that governance structures are chosen to minimize transaction costs under bounded rationality and opportunism. Separation has costs: coordination overhead, information loss between generator and verifier, slower iteration, the expense of maintaining two independent processes rather than one. In most organizational settings, the generator and evaluator are co-located precisely because the transaction cost of separating them exceeds the governance cost of co-location. This treatise does not argue that separation is unconditionally optimal. It argues that separation is optimal when the adversarial gradient (P1) is the dominant failure mode: when the cost of a generator gaming the evaluator exceeds the cost of maintaining structural separation. The boundary condition is: when stakes are low, when the generator is trusted, or when the speed premium of co-located iteration exceeds the governance benefit of separation, co-location may dominate. The treatise's empirical base is drawn entirely from settings where the adversarial gradient dominates (high-consequence reasoning under sustained optimization pressure), and the principles should be read as conditional on that regime rather than as universal prescriptions. A reader operating in a regime where the generator is not adversarial (where the primary failure mode is slow iteration rather than gaming) should weight the transaction cost of separation against the governance benefit before adopting Principle~I. A further Williamsonian axis the boundary condition must eventually address is asset specificity: a generator whose outputs are highly specific to a shared context produces co-location benefits that Williamson would classify as hierarchy-efficient even when opportunism is present. The current boundary condition names stakes, trust, and speed but omits specialization, and that omission makes the condition necessary but not sufficient for the governance choice it describes.

Separation must extend to the apparatus's own configuration decisions. A rubric author who selects reviewer personas based on knowledge of the ground truth introduces an oracle-lite contamination through the configuration layer. The defense is zero-oracle apparatus configuration: reviewer persona selection must be driven by observed failure signal (e.g.\ failure families from a latent-distance record), not by the operator's domain knowledge. Any configuration decision the generator cannot replicate from observable evidence alone is a candidate contamination channel.

\subsection{Principle II: Statelessness Without Social Cost}

The second principle is that each verification pass must be performed without memory of, or investment in, previous passes. This principle is the key to the system's structural difference from human adversarial review.

A human adversarial reviewer has identity, ego, and social reputation invested in their previous judgments. They cannot fully bracket the desire to be consistent with their last review, to please the principal, or to avoid the cost of admitting error mid-meeting. The cost of attacking one's own argument (or even of attacking an argument one has previously defended) is paid in social currency. The cost is a property of continuous selves with social histories, not a weakness of individual humans.

The verification process this treatise is about pays no social currency. Each pass is a fresh process with no memory of the previous pass's verdict and no investment in the previous pass's framing. The attack is \emph{zero-trust} in a way humans cannot be even when trying, not because humans are weak but because they are continuous selves. The property is structural, not motivational.

The organizational consequence is that verification can be made cheap to re-run. A human reviewer who has said an argument is sound cannot, without cost, say the next week that the same argument is unsound: the reversal costs them status. A stateless verification process can produce the reversal without cost, which is what makes it possible to apply the operations of Chapter~1 at the frequency the operations require. A once-a-quarter verification pass is a ceremony; a fifty-times-a-day verification pass is a discipline. Only the stateless version can run at the frequency the discipline requires.

The principle applies to human reviewers too, but only indirectly: a human reviewer who can be given structural permission to reverse a prior verdict without paying status cost can approximate the property. The case-method classroom is the clearest institutional simulation of this: students in good case discussions reverse positions multiple times per class because the norm of the classroom removes the social cost of doing so. The fact that the case method is the only widespread institutional simulation of this property is evidence that the property is rare and valuable.

A standing objection to this principle is that statelessness opens a slow-poison attack: an intelligent generator that knows each pass has no memory can execute incremental drifts that satisfy every single pass individually and only constitute a failure when viewed chronologically. The objection is real and the answer is that statelessness is per-pass, not per-run. The detection of chronological drift is the job of Principle~IV (cheap repetition) composed with Principle~VI (holdout surfaces authored outside the candidate's claim region): drift shows up as a residual on a forward-observable region the generator did not shape, sampled across runs. The stateless pass is the atomic unit; the across-pass residual on an authored holdout is the slow-poison detector. Statelessness without IV+VI would indeed be a vulnerability; statelessness inside the full principle set is what makes drift detectable without giving any single pass the reviewer's memory bias.

\subsection{Principle III: Typed Operations Bound to Deterministic Checks}

The third principle is that the operations of Chapter~1 must be bound to checks a third party can evaluate without trusting the operator. A verification practice in which the operations are performed at the level of intuition, without externalized checks, cannot produce repeatable results and cannot be audited.

The principle is easiest to see by contrast. Consider a human reviewer who says ``this argument has a load-bearing claim problem.'' The statement is a judgment. It cannot be directly checked. Another reviewer can agree or disagree, but the question of whether the load-bearing claim has in fact been correctly identified is resolved by appeal to intuition and consensus, not by a procedure a third party can run.

Now consider the same reviewer working within a system where ``load-bearing claim'' is a named slot in a structured argument representation (where the argument has been parsed into a graph of claims, where each claim has been tagged with its dependencies, and where the ``load-bearing'' designation corresponds to a specific graph-theoretic property: the claim whose removal produces the largest downstream collapse). In that system, the load-bearing claim can be identified by a procedure. The procedure can be wrong, but it can be wrong in a way that is checkable. Another reviewer can run the procedure on the same argument and see whether they get the same answer. Disagreements are resolvable by examining the procedure, not by appealing to authority.

Operations bound to deterministic checks are what make the decomposition teachable and auditable. A verification practice in which every operation is a judgment cannot be improved, because its failures cannot be localized. A practice in which each operation is bound to a check can be improved, because a failure of the check is a specific thing that can be debugged.

This is the principle that most separates the decomposition in Chapter~1 from incumbent critical-thinking curricula. Incumbent curricula name the operations (if they name them at all) as dispositions or habits of mind. The decomposition here names them as typed slots in a process that can be externally checked. The difference is structural: it separates a practice that can accumulate from one that cannot.

\subsection{Principle IV: Cheap Repetition Without Exhaustion}

The fourth principle is that the verification process must be cheap enough to run many times on the same argument, and the process must not degrade with repetition. A verification practice that is expensive per pass, or that becomes less rigorous with each subsequent pass, cannot catch the failure modes of Section~1.3 reliably, because those failure modes are often detectable only after repeated pressure.

Human reviewers degrade with repetition for well-understood reasons: fatigue, boredom, confirmation bias, the desire for closure, the cognitive cost of sustained skepticism, and the social cost of continuing to attack an argument that has already been attacked. A human reviewer's first pass is sharper than their fifth pass. A human reviewer's fiftieth pass is a rubber stamp.

A decomposed verification process built on the principles above does not degrade in this way. Each pass is stateless (Principle~II), typed (Principle~III), and structurally separated from the generator (Principle~I). The fiftieth pass is as ruthless as the first, because the fiftieth pass is in fact a fresh process that knows nothing of the previous forty-nine except for the typed state they left behind.

The consequence of this property is that the verification practice can be applied at frequencies that would be impossible for humans. An argument that has been pressured fifty times by a stateless verification process has been pressured in a way that no amount of human review can replicate. The failure modes that survive ordinary human review (the ones the field manual catalogues) are exactly the ones that only become visible under sustained pressure. A one-pass review is structurally incapable of catching them. A fifty-pass review is structurally capable of catching them if and only if the fifty passes are each as sharp as the first.

This is the property that most distinguishes the verification practice this treatise is about from human adversarial review. Humans can do the operations of Chapter~1 well; they cannot do them at the frequency required to catch the pathologies of Section~1.3 reliably. The frequency is the principle; the operations are only useful if they can be run at it.

\subsection{Principle V: Pre-registration of the Killer Question}

The fifth principle is that the verification criterion (the question that, if answered wrongly, would falsify the argument) must be fixed \emph{before} the argument is evaluated against it. A verification practice that allows the criterion to be chosen after the argument has been generated is a practice in which the generator can steer the criterion toward claims it already satisfies. This is the standard pathology of ex-post evaluation.

The principle is borrowed from experimental science, where pre-registration of the hypothesis and the test is the standard defense against post-hoc rationalization. It applies to epistemic verification for the identical reason: once an argument exists, the space of possible criteria that would score it favorably is large enough that a sufficiently motivated evaluator can always find one. The only defense against this is to commit to the criterion before the argument is produced.

In practice, this means that a verification pass must begin with the question \emph{what would falsify the load-bearing claim here?} and must record the answer before the argument is attacked. The answer is the killer question in its operational form. The attack that follows is the attempt to produce the falsifying observation. If the falsifying observation is not produced, the argument survives; if it is, the argument does not. The sequence matters. A verification practice that produces the killer question after the attack has already found its target is laundering the same pathology the practice is designed to catch.

\subsection{Principle VI: Holdout Test Surfaces Authored Outside the Candidate's Claim Region}

The sixth principle is that the test surface against which a candidate is evaluated must be authored outside the region the candidate's own claim inhabits. Pre-registration (Principle~V) fixes the criterion before the argument exists; this principle fixes the \emph{domain of observation} so that the argument cannot steer the verifier into a region where its claim is locally well-behaved.

The distinction matters because a candidate argument optimized against its own fit window is almost always persuasive inside that window. The Ptolemaic precedent is the clearest case: a deferent-and-epicycle model of planetary motion is arbitrarily accurate on the orbits it was fit to, and its failure only becomes visible when the test is performed on orbits it was not fit to. A parsimonious model that scores perfectly on the fit window is not more trustworthy than a messy one; it is often \emph{less} trustworthy, because parsimony is what makes the fit-window score persuasive, and persuasion is exactly the channel the failure travels through.

The principle is therefore: whenever a verification pass is evaluating a claim with asymptotic structure (a scaling law, a limit behavior, a forecast that extends beyond the observed window, a thesis that generalizes from a local case) the pass must generate test cases from a region the candidate did not see, did not cite, and cannot have shaped. The test region must be authored by a process independent of the candidate. If the candidate authored any part of the test surface, or if the test surface was derived from the candidate's own output, the test has been structurally compromised in the same way Principle~I's violated separation compromises the verifier.

The practical consequence is that a verification process that scores a candidate at 100 on its training window and never tests a held-out region has not performed verification in the sense of this principle. It has performed fit. The two are indistinguishable to an observer who only sees the in-window score, and this indistinguishability is the central reason why verification that relies only on in-window metrics is structurally incapable of catching a class of failures that are specifically visible outside the window. A candidate's ability to pass a holdout it did not author is the move that distinguishes a surviving hypothesis from an elegant locally-correct surrogate.

This principle is the third leg of the verification apparatus, alongside separation (Principle~I) and pre-registration (Principle~V), and its omission from an earlier framing of this treatise is recorded as a load-bearing correction: a two-principle description (separate the generator from the verifier, pre-register the criterion) is insufficient, because it leaves intact the move by which a candidate authors its own test surface. Without the holdout-region requirement, the verifier can still be steered, and the apparatus collapses to something indistinguishable from a better fitter with a falsifier bolted on.

The principle is also the operational form of a distinction drawn several decades earlier in the philosophy of science. \citet{lakatos1978} distinguished \emph{progressive} from \emph{degenerating} research programmes by asking whether the programme, when confronted with an anomaly, responded by adding auxiliary hypotheses chosen to fit the known data (degenerating) or by predicting a novel observable outside the region the programme was built to fit and then surviving the test of that observable (progressive). Principle~VI is the structural requirement that converts the Lakatosian distinction into a runnable verification step: a candidate's survival on a test surface authored outside its own claim region is, by construction, a progressive move in Lakatos's sense; a candidate's survival on a test surface it authored or shaped is, by construction, a degenerating move: the auxiliary was built to the data, and the ``survival'' is a trivial consequence of the authorship. A verification apparatus that cannot distinguish these two is not distinguishing a progressive programme from a degenerating one, which is the central distinction Lakatos argued any honest methodology of science must be able to make. This treatise's claim is that Principle~VI operationalizes that distinction for contemporary LLM-driven verification: the forward-observable test surface, authored by a process independent of the candidate, is the progressive-vs-degenerating discriminator that fits the current substrate. Lakatos originally framed this distinction temporally: a progressive programme predicts a novel fact that is later confirmed. In the context of computational verification, the temporal requirement translates into a spatial one: the test surface must be held out from the candidate's fitting process. A candidate surviving a spatial holdout it did not author is functionally analogous to a theory surviving a temporal prediction it could not have retrofitted, provided the holdout author is independent of the generator and does not share its priors. The independence condition is what makes the spatial translation legitimate; without it, the holdout author can unconsciously construct a test surface the generator is already equipped to pass.

\subsection{Principle VII: Asymptotic Rather Than Absolute Scoring Standards}

The seventh principle governs how verification outcomes are reported once the apparatus of Principles~I--VI has been applied. The verification standard must be asymptotic (unreachable in the limit) rather than absolute. A standard an argument can satisfy completely is a standard an argument will be optimized to satisfy, at which point the standard ceases to discriminate between good arguments and arguments that are good at passing the standard.

This is the Goodhart property. Any measure that becomes a target ceases to be a good measure, because the optimization pressure of targeting it destroys the property that made it measurable in the first place. A verification standard that can be fully satisfied (where there exists a score of 100) will produce arguments that achieve the 100 without being good, because the optimization surface is closed.

The defense is to set the standard so that full satisfaction is structurally impossible. Not because the standard is arbitrary, but because the thing the standard is trying to measure (the soundness of an argument under arbitrary future pressure) is genuinely open-ended. An argument can be pressured indefinitely, and the evidence that it has survived this pressure can always be strengthened. A standard that recognizes this by remaining open (by always leaving a next question to ask) preserves its discriminative power indefinitely. A ceiling below the nominal maximum is one way to operationalize this; other implementations are possible and the principle is portable in a way any particular ceiling value is not.

\subsection{A Standing Reservation: The Library Is Itself a Goodhart Target}

The seven principles above specify an inspection architecture. They do not, by themselves, defend the architecture against the failure mode the architecture is designed to detect. A gate library is a specification, and any specification an optimizer can read is a specification the optimizer will eventually learn to satisfy without doing the underlying work. The ten operations of Chapter~1 and the seven principles of Chapter~2 are therefore themselves targets the next generation of generators will optimize against. A candidate that passes every gate in the library without producing a truly novel contribution is possible in principle and observed in practice.

The partial answer is Principle~VII: an asymptotic scoring standard that cannot be fully satisfied keeps the surface open and the library's effective cardinality growing. The more honest answer is that this is an arms race, not a fixed defense. The gate library has to turn over (new compositions must be added, old ones retired, and the turnover has to be authored outside the region recent candidates have shaped) at a rate that matches the generator's adaptation rate. A static library of fixed cardinality will eventually be gamed no matter how carefully its gates are designed; a library that turns over faster than the generator can reverse-engineer its composition retains its discriminative power. The operator's role in running this turnover is one of the residual operations named in Chapter~3, and is not part of the decomposition.

The reservation is load-bearing because a reader who takes the seven principles as a finished architecture will miss the thing they are embedded in: a process of continuous library extension, performed by a party the generator cannot predict, against a moving target. The treatise names the principles because they are the portable claims. It names this reservation because the principles are not self-enforcing and pretending otherwise is the specific failure mode the principles are designed to catch.

This turnover mechanism also serves as the apparatus's defense against distribution shift. The decomposed operations verify an argument against a static holdout (Principle~VI) at design time. The ledger (\S2.9) detects when the environment has shifted at runtime, forcing the operator to author new gates when old ones cease to discriminate. Distribution shift detection is not an eleventh operation of epistemic verification; it is a monitoring function of the ledger, and the distinction matters because conflating the two would expand the decomposition's scope beyond argument-level verification into environment-level surveillance.

\subsection{The Static Grammar as Falsification Guarantee}

The arms race described in \S2.8 generates a natural counterproposal: if the gate library must turn over, why not eliminate the primitive library entirely and allow the generator to discover in an unconstrained representation space? The static vocabulary of symbolic primitives (power laws, exponentials, rational forms) looks, from this angle, like an artificial constraint on discovery. The counterproposal reaches its most sophisticated form in symbolic distillation pipelines: let a neural network discover in its own high-dimensional latent space, then translate the result into human-readable algebra. On this view, the grammar is a display layer, not a discovery boundary.

This framing has the relationship backwards. The static grammar is not a constraint on discovery. It is the guarantee that the farther-tail gate's verdict is about the proposed law rather than a translation of it.

The argument is structural. Suppose the generator proposes an expression drawn from a fixed symbolic vocabulary~$G$. The farther-tail gate evaluates whether that expression correctly predicts observations at points far outside the evidence grid. If it fails, the failure is attributable to the expression itself~--- the symbolic form is structurally wrong for the domain, and the gate's verdict is about the claim, not about the mechanism that produced it. The static grammar guarantees this traceability. The gate tests the law as stated, not a surrogate for the law.

Now suppose the generator discovers in an unconstrained neural representation and a symbolic distillation step translates the result into algebra. The farther-tail gate now tests the output of the translation step. A failure at the gate can mean one of two things: the underlying discovery is wrong, or the translation step lost information that the discovery contained. These two failure modes are indistinguishable without access to the latent representation~--- and if that representation were accessible, the symbolic distillation step would not be needed. The gate's verdict is no longer about the law. It is about the law as the distillation step expressed it. This is an unauditable trust gap, not an expansion of the discovery space.

The practical version of this gap appeared in the gp023 crucial experiments on a two-variable transcendental substrate. The visible evidence grid was consistent with several structurally distinct forms~--- Wien approximation, Weibull-type exponential, and the true ground-truth law~--- because none of these forms diverges substantially within the observed regime. The farther-tail gate at holdout points an order of magnitude beyond the evidence boundary discriminates them: the approximating forms overpredict the ground truth by two to four orders of magnitude in the tail; the correct structural law does not. A generator operating in the static symbolic grammar cannot express the correct form if the requisite denominator structure is outside the grammar, and the gate correctly catches this: the expression fails at the discriminating regime. The gate's verdict is diagnostic~--- it names the structural deficiency. A generator discovering in a neural latent space would produce a curve that fits the visible grid and fails the farther-tail gate for the same structural reason, but the failure would be attributed to the distillation step rather than to the underlying discovery, making the diagnosis unavailable.

The deeper point is that a dynamic grammar~--- one that expands to accommodate any expression the generator discovers~--- is structurally equivalent to a neural network under a different name. A neural network is a dynamic grammar with a continuous parameterization. Its fitting power is high precisely because it has no primitive constraint: it can approximate any smooth function on the observed evidence to arbitrary precision. The farther-tail gate cannot distinguish between ``the network discovered the correct structural law and expressed it faithfully'' and ``the network found a high-fidelity interpolant that fails to extrapolate.'' There is no symbolic structure to interrogate. There is only a prediction at a new point, which a sufficiently flexible approximator can produce correctly by local interpolation even when the underlying form is wrong.

The static grammar is therefore not an obstacle to be engineered around. It is the mechanism by which the farther-tail gate's verdict remains about the proposed law rather than about the approximation of the proposed law. The operator-guided grammar expansion described in \S2.8~--- introducing new structural primitives when the Falsification Suite identifies the current library as exhausted~--- is the correct mechanism for breaking the score ceiling not because it is the only way to add new symbolic forms, but because it is the only way to add them while preserving the gate's epistemic traceability. Each new primitive is a falsifiable structural commitment, testable by the farther-tail gate on the next run. The commitment is what makes the gate's verdict informative rather than definitional.

\subsubsection*{Grammar Expressiveness as a Measurable Quantity}

A corollary of the static grammar's role as falsification guarantee is that grammar expressiveness is itself a measurable property of the verification apparatus, not merely a design choice. In continuous-optimization mode, a stagnation ceiling is ambiguous: it could reflect optimizer pathology (local minima, initialization sensitivity) or genuine grammar insufficiency. The two failure modes require different remedies (more restarts vs.\ primitive injection), and the apparatus must provide the instrument that separates them~--- multi-start restarts with residual spread analysis, where high spread across restarts indicates optimizer pathology and low spread at a persistently bad residual indicates grammar ceiling.

In discrete evaluation mode~--- when scoring is exact-match over a countable output space and the optimizer has no continuous surface to descend~--- the ambiguity collapses. A stagnation ceiling at maximum error is a classification-complete event: optimizer pathology is ruled out by architecture (there is no optimizer foothold), and grammar insufficiency is the only remaining interpretation. The ceiling is no longer a signal that requires disambiguation; it is a direct measurement of a property of the grammar on the given substrate. A grammar that stagnates on a discrete target is one that does not contain a correct expression for that target, and that is a fact about the grammar, not about the search procedure.

This distinction matters for the interpretation of empirical results. When an apparatus confirms grammar insufficiency on a discrete substrate, the confirmed claim is that the grammar lacks the required form~--- not that the same grammar, applied without retrieval constraints, would also fail to find the correct form by label lookup. Grammar insufficiency and retrieval blocking are two distinct mechanisms that can each produce score ceilings and must be evidenced separately. An experiment that blocks label retrieval (via a named-import gate and score-zeroing) and an experiment that operates without such a gate can produce different outcomes on the same grammar and the same substrate; the former tests structural articulation capacity, the latter tests whether the grammar even contains the correct form. Conflating these two results produces the strongest-sounding claim with the weakest evidence. The conservative statement is always: confirm the claim the experiment actually controls, and name what would need to differ for the stronger claim to be warranted.

\subsection{The Epistemological Ledger: Operationalizing Belief Under Optimization Pressure}

The seven principles above specify what an inspection architecture must look like at any given moment. They do not specify how the architecture tracks what it has learned across time: how a hypothesis becomes an experiment, an experiment becomes a finding, and a finding becomes actionable belief. This section describes a three-tier recording discipline (the epistemological ledger) that operationalizes Principles V (pre-registration) and VI (holdout surfaces) at the project-governance layer and provides the turnover mechanism \S2.8 identifies as necessary for the library to remain discriminative.

The ledger is offered as a transferable methodology, not as a description of one system's bookkeeping. A reader building a different verification system for a different purpose needs the same structural answer to the same structural question: how do you prevent your own beliefs from drifting under the optimization pressure your system generates? The ledger is one answer. It is the answer this project converged on after discovering, through repeated incident, that each of the seven principles can be satisfied locally while the project's aggregate confidence drifts globally.

\subsubsection*{The Three Tiers}

\textbf{Tier 1: Hypothesis registration.} Before a discriminating experiment runs, the operator records the hypothesis, the eigenquestion, the falsifiable claim, the test that would discriminate, and the success criterion. The record is sealed before the first data point is observed. This is Principle~V (pre-registration) applied not to a single verification pass but to the project's evolving belief about what is true. A hypothesis that is registered after the experiment has already produced its result is flagged as backfilled and loses evidential standing: it may be recorded for completeness, but it cannot update the project's confidence on the claim it names.

The discipline matters because optimization pressure acts on the operator's posterior as well as on the generator's output. An operator who has watched forty iterations of a verification loop will, without the ledger, selectively remember the iterations that confirmed their developing intuition and discount the ones that did not. Pre-registration is the defense against this drift. It is expensive (formulating a falsifiable hypothesis before the run is harder than narrating the result after) and it is non-negotiable for the same reason Principle~V is non-negotiable: the alternative is a system whose beliefs are shaped by the same ex-post rationalization the system exists to detect.

\textbf{Tier 2: Experiment closure.} When an experiment ends (whether by completing its planned iterations, by operator stoppage, or by apparatus failure) the result is recorded as a closed fact. The record captures what ran, what the actual result was (including null results), what changed in the project's understanding, and where the source artifacts live. Not every experiment changes what the project believes; many produce null results, apparatus calibrations, or replications of known patterns. These are recorded with the same discipline as positive results, because a project that only records its successes is a project whose ledger is itself a Goodhart target.

The closure discipline has a specific test: if a cold agent (one with no memory of the conversation that produced the experiment) opens the project tomorrow, can it find what ran, what was learned, and what to do next without reading chat history? If not, the closure is incomplete. This test operationalizes Principle~II (statelessness) at the governance layer: the ledger must be readable by a stateless process, not only by the operator who ran the experiment.

\textbf{Tier 3: Finding promotion.} A finding is a belief update. Not every experiment produces one. A finding is recorded only when the experiment changes what the project should believe or build next, and the finding must cite the hypothesis it tested and the experiment that produced the evidence.

The promotion gate is the load-bearing structural element: a finding can be marked \emph{active} (available to inform subsequent decisions) only if the pattern it names has been observed at least twice across independent runs. A pattern observed once is recorded at \emph{note} status: acknowledged, tracked, but not actionable. The two-strike rule is the ledger's defense against the single-instance illusion: the tendency for a vivid one-off result to become doctrine because it confirmed a prior the operator was already holding.

The rule has two exceptions, both operator-controlled: an approved verifier experiment that would produce the second instance if the pattern is real, or an operator-grade decision explicitly marked as independent of evidence. Both exceptions are recorded in the ledger so that a future reader can see which findings were promoted on evidence and which on operator judgment. The distinction matters because operator-promoted findings are the ones most likely to be wrong, and the ledger's job is to make that likelihood visible rather than to hide it.

\subsubsection*{Why the Ledger Is Not Bookkeeping}

A reader who treats the three tiers as project management will miss the structural claim. The ledger is not a record of what happened. It is a \emph{constraint on what the project is allowed to believe}. Each tier enforces a specific principle:

\begin{itemize}
\item Tier~1 enforces Principle~V: you cannot claim a result you did not pre-register.
\item Tier~2 enforces Principle~II: the record must be readable without the operator's memory.
\item Tier~3 enforces Principle~VII: belief is asymptotic (a single observation cannot close the question) and the two-strike rule is the operational form of that asymptote.
\end{itemize}

The three tiers compose into a defense against a failure mode no individual principle catches: \emph{slow belief drift under optimization pressure}. An operator running a verification system generates optimization pressure on their own posterior. Every iteration produces a result. Every result nudges the operator's confidence. Without the ledger, the operator's beliefs are updated continuously by a stream of observations whose evidential weight varies enormously but whose psychological salience is roughly uniform. The ledger forces the operator to distinguish between observations that meet the evidential bar (promoted findings) and observations that do not (notes), and to act only on the former.

The ledger also provides the turnover mechanism \S2.8 identifies as necessary. The standing reservation warns that a static gate library will eventually be gamed. The ledger tracks which gates have been tested, which failure modes have been observed, and which hypotheses have been falsified. This record is the input to library turnover: new gates are authored in response to observed failure modes (Tier~2), and the decision to add them is justified by promoted findings (Tier~3) rather than by the operator's intuition. The ledger does not solve the arms race. It makes the arms race visible and auditable, which is the most that any governance layer can do.

\subsubsection*{A Live Instance: The Misfile}

The pathologies of \S1.3 are not confined to strategic arguments under evaluation. They occur in the apparatus's own metadata. A concrete instance, discovered during the development of the system this treatise draws on, illustrates both the pathology and the ledger's role in catching it.

A library of mathematical corrector forms (candidate functions the verification apparatus tests against observed data) contained a metadata field classifying each form as \emph{smooth} or \emph{non-smooth}. The classification was used by a downstream filter to narrow the candidate pool before a deterministic sweep. The function $\text{round}(k \cdot v)$ (which produces a step function with jump discontinuities at every half-integer) was classified as \emph{smooth} by the programmer who entered it, because \texttt{round()} reads as a standard mathematical operation to human intuition.

The misclassification was Pattern~7 (The Misfile) in its purest form: a fact placed in the wrong category early in the pipeline, with every downstream inference inheriting the error silently. The verification apparatus worked flawlessly: it correctly identified the data as step-shaped, correctly queried the library for non-smooth forms, and correctly reported that the best surviving candidate was $\text{floor}(k \cdot v)$ with a fitted parameter of 0.097. The answer was wrong because the dictionary was wrong, not because the reasoning was wrong. No code crashed. No error was raised. The apparatus confidently guided the generator toward a form that did not match the data, and the generator (unable to reconcile the guidance with what it could see) produced increasingly desperate alternatives that scored zero.

The fix was an automated metadata validator: a unit test that empirically measures each library form's discontinuity and monotonicity across a sample range and asserts the declared metadata matches. The validator caught eleven additional misclassifications beyond the one that triggered the investigation. The ledger's role was to ensure the incident was recorded as a finding (not merely fixed as a bug), with the structural lesson (\emph{do not trust human intuition about the calculus properties of the apparatus's own primitives}) promoted to active status and available to inform future library extensions.

The incident is worth recording in a treatise about epistemic verification because it demonstrates that the apparatus is not exempt from its own pathology catalogue. A verification system that catches Pattern~7 in the arguments it evaluates but not in its own metadata is a system whose inspection principle (\S A.2, Condition~5) has a hole at the self-referential layer. The automated validator (a deterministic check on the apparatus's own dictionary) is the same structural move as Principle~III (typed operations bound to deterministic checks), applied inward.

%% ──────────────────────────────────────────────────────────────────────────────
\subsection{From Verification to Discovery: Three Adversarial Gates}

The empirical mining and the self-correcting audit above concern verification: the apparatus testing claims that have already been generated. A separate question is whether the operations of Chapter~1, composed into a recursive loop with deterministic gates, can perform the upstream activity of \emph{discovering} a functional law from data rather than merely \emph{verifying} one that a human or an LLM has proposed. The distinction matters because a system that only verifies is an instrument; a system that also discovers is an engine. The treatise's scope commitment (\S 1.2) does not claim that the apparatus is a discovery engine, but the empirical record now contains three pre-registered synthetic tests that bear on whether it could become one, and honest reporting requires that they be disclosed.

The three tests were designed to catch specific failure modes that would disqualify a discovery claim even if holdout accuracy were high. Each test uses a synthetic substrate---a ground-truth law authored by the operator, sealed before the run, with data points visible to the mutator but the generating law and its constants hidden.

\paragraph{Test 1: Retrieval versus Discovery.} The substrate uses the functional form $\alpha(d) = C_1 / (d + C_2)$---structurally identical to published scaling-law relationships---but with constants $C_1 = 3.714$ and $C_2 = 0.892$ that do not appear in any published work. The known literature values are $C = 2$ or $C = 4$ with zero offset. An anti-retrieval gate checks whether the discovered constants are within 5\% of any known value; if so, the apparatus has retrieved from parametric memory rather than discovered from data. Result: the apparatus recovered $a = 3.703$, $b = 0.886$ (error ${<}0.3\%$ on both constants), with holdout mean relative error 0.8\% on five withheld points and 3.9\% on four extrapolation points beyond the visible range. The anti-retrieval gate confirmed the constants are distinct from all published values. Score: 82/100, with the residual 18 points attributable to the judge's correct observation that the denominator offset is empirically fitted without mechanistic grounding.

\paragraph{Test 2: Asymptotic Discipline.} The substrate uses a two-term decay law (exponential plus power-law) that approaches zero from above as $d \to \infty$. A polynomial approximation, no matter how accurate on the visible range, will extrapolate to negative values at large~$d$. An asymptotic-wall gate probes the discovered law at $d \in \{100, 150, 200\}$ and rejects any prediction outside $[0, 1]$. Result: the apparatus discovered a form that passes all holdout gates (MRE 1.3\%) and stays within bounds at all three probe points. A degree-3 polynomial fitted to the same visible data extrapolates to $-32$ at $d = 100$. Score: 90/100.

\paragraph{Test 3: Parsimony Resistance.} The substrate uses a deliberately complex ground truth with $K = 10$ parameters, including an oscillatory term superimposed on a decay envelope. A five-parameter monotonic model fits the visible data acceptably (visible MRE ${\approx}7.6\%$) but fails holdout (MRE ${\approx}20\%$) because it smooths over the oscillatory structure. The test checks whether the apparatus forces parsimony---accepting the simpler, wrong model because BIC rewards fewer parameters---or accepts the messy truth when holdout demands it. Result: the apparatus accepted a higher-parameter form that captures the oscillatory structure. Holdout MRE: 4.9\%. Score: 90/100.

\paragraph{Cross-mutator replication.} Each test was independently replicated with a second mutator model from a different provider family: OpenAI o3 (reasoning model) and Anthropic Claude~Opus or Claude~Sonnet (general model). Results: GP-159 scores 90 (o3) and 90 (Claude~Sonnet). GP-160 scores 90 (o3) and 82 (Claude~Sonnet). GP-161 scores 90 (o3) and 81 (Claude~Opus). All three tests pass across both mutator families, confirming the result is structural---a property of the apparatus and the gate architecture---not an artifact of a single model's training distribution.

\paragraph{Scope of the claim.} The three tests, replicated across mutator families, support a scoped methodological finding: the apparatus has the raw capability to enable scientific discovery. It avoids three specific failure modes (retrieval, extrapolation breakdown, and MDL Goodharting) that would disqualify any discovery claim, and it does so regardless of which frontier model performs the generation step. This does \emph{not} establish that the apparatus discovers laws on real-world scientific substrates where the ground truth is unknown. That claim requires at least one real-world substrate where the apparatus proposes a law that survives independent empirical validation, and one substrate where neither the operator nor the apparatus has access to the ground truth at design time. The analogy is to instrumentation: the three tests prove the telescope works; they do not prove the telescope has found a new planet. The synthetic triad is a necessary precondition, not a sufficient demonstration.

%% ──────────────────────────────────────────────────────────────────────────────
\section{Chapter 3: The Residual: What Did Not Systematize}

The decomposition in Chapter~1 is real, and the principles in Chapter~2 are load-bearing for why it works. If the treatise stopped there, it would be a Taylorist document in the bad sense: a decomposition held up as complete without honest accounting of what the decomposition failed to capture. This chapter is the counterweight. It names the residual.

The residual is the portion of epistemic verification that has not, so far, decomposed into operations the stateless, typed, deterministic process of Chapters~1 and~2 can perform. This chapter claims that the residual is structurally different from the operations that did decompose, and offers a reason (grounded in \citeauthor{peirce1878}'s distinction between abduction and the other two inference types) to expect the residual to resist decomposition longer than the craftsmen's residual resisted Taylor's. But a treatise whose central argument is that guilds misjudge the boundaries of the ineffable must hold its own boundary claims to the same standard. The honest framing is: the three operations below have not decomposed \emph{on this system, under this operator, over four months}. Whether they are permanently irreducible or merely not yet decomposed is an empirical question this treatise cannot answer from the inside. The Peircean argument provides structural reasons to expect resistance; it does not provide a proof of impossibility. If a future system decomposes eigenquestion selection into a stateless, typed operation without loss of the property that makes the operation load-bearing, this chapter's claim is refuted and should be updated rather than defended.

The distinction between the craftsmen's residual and this one is worth naming explicitly, because the treatise's own Taylor analogy creates the obvious objection. The craftsmen in 1911 insisted their residual was irreducible, and Taylor showed it was not. Why should this residual be different? The answer offered here (and it is offered as a hypothesis, not a proof) is that the craftsmen's residual was sensorimotor: the feel of the metal, the angle of the chisel, the weight of the pour. These are operations on physical substrates that were not yet instrumented. Once instrumented, they decomposed. The residual named below is abductive: the generation of a frame before there is a procedure to generate it. \citeauthor{peirce1878}'s central claim is that abduction is not reducible to deduction or induction by composition: it is a different logical type, not a different difficulty level. If Peirce is right, the residual resists decomposition for a reason the craftsmen's did not: not because no one has yet built the instrument, but because the operation is of a type the instrument cannot perform. If Peirce is wrong (if abduction turns out to be induction over a latent space that has not yet been featurized) then the residual narrows, and the treatise's boundary moves. Either outcome is informative. Neither outcome damages the decomposition of Chapters~1 and~2, which stands or falls on its own evidence regardless of where the residual boundary eventually settles.

Three operations have resisted decomposition. They are named and described below.

Before listing them it is worth naming the psychological signature they share, because the boundary this chapter draws becomes sharper once the shape of the residual is visible as a shape rather than as a list. In the vocabulary of self-determination theory \citep{ryan2000}, an activity is experienced as \emph{intrinsically motivated} (chosen from self-endorsed volition rather than performed under external contingency) when it satisfies three basic needs: autonomy (the actor experiences the action as self-originated), competence (the actor experiences the action as matched to real capability), and relatedness (the action has standing in a community whose judgment the actor is accountable to). The decomposed operations of Chapter~1 do not satisfy these needs and do not need to; they are typed, stateless passes whose credibility comes from structural independence rather than from the psychological state of whatever is performing them. The three operations below are different in exactly this respect. Each one's value is partly constituted by being performed under autonomous commitment: by an agent who chose this frame, in this moment, for this reason, with standing in a community that will evaluate the choice. The claim is not that the operations are mysterious. They live in the zone self-determination theory names explicitly, and the residual is therefore bounded in a way the incumbent vocabulary of ``judgment'' is not: it is the zone in which autonomous, competence-matched, community-accountable commitment is load-bearing, and the structured-input / structured-output mode of Chapters~1 and~2 is architecturally unsuited to producing it. The three operations named below are, read in this light, the three places where the SDT zone enters the verification practice.

\subsection{The Selection of the Eigenquestion Itself}

The clearest instance of the residual this chapter names happened on 2026-04-14, and it is worth leading with because the abstract claim lands differently once the live work is visible.

While preparing to seal a stress-test sandbox for the verification engine, the operator declined a pre-commit identifiability check that had returned a clean pass. The check was a bootstrap-under-noise test: the optimizer recovered the declared ground-truth parameters consistently across noise realizations from a fixed starting point, and by the decomposed rules of the apparatus this was sufficient for the sealing gate. The operator ran a second check anyway (an adversarial multi-start fit from clean data) on no stronger warrant than a general commitment to distrust green lights that arrive too easily. The second check failed: two of the six declared parameters were recovered with a 70\% error and a 1:1 ratio to each other identical to their ratio in the ground truth, which on inspection revealed that the two parameters entered the functional form only through their quotient, making the declared six-parameter family structurally rank-five. The apparatus caught the specification flaw, but only because the operator pointed the apparatus at a question it was not required to ask.

The decomposed apparatus is what made the second check runnable; the residual is what decided to run it. That decision is what \S3.1 names. It is not a failure of the decomposition to reach further (the apparatus performed exactly as specified once pointed) and it is not evidence that the operator is heroically indispensable. It is evidence that the decomposed operations cannot produce the commitment to run a check the rules did not mandate, and the operator under the discipline of distrusting green lights can.

Now the abstract claim. The first operation the decomposition has not captured is \emph{which question should be eigenquestion-identified in the first place}. Operation~1 of Chapter~1 takes as input a set of claims and selects the one whose answer determines the rest. The input (the set of claims) is itself the output of a prior decision about what argument to attack, and that prior decision has not systematized. When the verification process is handed an argument, it can identify the eigenquestion within that argument reliably. When the verification process is asked to decide \emph{which argument} to verify, \emph{which reframing of the underlying problem} to work within, or \emph{which check to run that the rules did not mandate}, it does not have a procedure that produces an answer. The operation is performed, in the existing system, by the human operator.

The logical character of the operation can be stated precisely. \citet{peirce1878} distinguished three kinds of inference: \emph{deduction}, from rule and case to result; \emph{induction}, from case and result to rule; and \emph{abduction}, from result to the spontaneous generation of an explanatory hypothesis that, if true, would make the result a matter of course. The decomposed operations of Chapter~1 are deductive in Peirce's sense: they apply fixed rules to structured inputs. The generator the verifier attacks is inductive: a statistical substrate that finds rules from examples. The selection of the eigenquestion is neither. It is abduction: the operation by which a frame is proposed at all, before there is a rule to apply or a case to induce over. Peirce's central observation was that abduction is not reducible to the other two, and that a system architecturally limited to deduction and induction cannot produce it by composition. The eigenquestion-selection residual is therefore not a decomposition the apparatus has failed to reach yet; it is, on Peircean grounds, a category of inference the apparatus is architecturally unable to perform. The 2026-04-14 catch above is what this category looks like when it does one unit of load-bearing work in real time. The residual is not a gap in engineering. It is a gap in the logical type of operation the decomposed apparatus can contain.

The incumbent vocabulary calls this operation \emph{strategic judgment} or \emph{research taste} or \emph{knowing what to work on}. Those labels are doing real work. The operation is load-bearing, it is teachable only by apprenticeship, and it is the one place in this treatise where the guild defense has genuine force. It is not decomposable into stateless, typed passes. It is what stays human.

\subsection{Recognizing When to Reframe Rather Than Attack}

The second operation that has not decomposed is \emph{recognizing that the right move is a reframing rather than an attack}. The operations of Chapter~1 contain both basin search (Operation~6) and topological pivot recognition (Operation~3), and a competent verifier can perform each. What has not systematized is the decision of which of the two to perform on a given argument in a given state.

The decision is not arbitrary. There are signals (persistent stagnation at the current frame, recurring attack patterns that fail for the same underlying reason, the sense that a particular assumption is quietly doing all the work) that indicate a topological pivot is the correct move. The decomposed process can notice these signals when they are bright. When they are dim, the decomposed process tends to keep running basin searches, because basin search is the cheap default and the cost of a missed reframing is paid in delayed discovery rather than in visible failure.

In the existing system, the call between basin search and reframing is made, again, by the human operator. The operator reads the workspace, sees that the iterations have been producing the same failure signatures for too long, and forces a reframing. The decomposed process cannot produce this call reliably on its own. The signals it can read are a proper subset of the signals the operator reads; the residual is not yet specified precisely but is persistent across domains.

The incumbent vocabulary calls this \emph{seeing the bigger picture} or \emph{knowing when to step back}. Like the first residual operation, the label is doing real work. The operation is distinct from basin search and from topological pivot in the same way that selecting a problem is distinct from solving a problem: it is a meta-level commitment that the decomposed process has not been built to produce.

\subsection{The Social Dynamics of Live Pressure-Testing}

The third operation that has not decomposed is the performance of adversarial verification in a live, multi-party, social setting. The decomposed process operates on written arguments, at the operator's discretion, with no audience. Much of the practice this treatise is about (the diligence room, the board meeting, the case-method classroom, the peer review) is performed live, in front of an audience, under social stakes.

The operations of Chapter~1 are directly applicable to the content of the argument in these settings. The pathologies of \S1.3 are exactly the pathologies these settings produce. But the \emph{delivery} of the verification in a live setting (when to press and when to let pass, when to let the author save face and when to force them into contact with the weakness, when to escalate and when to retreat, when to speak and when to let the room do the work) is an operation the decomposed process cannot perform, because the decomposed process has no theory of the room.

The distinction worth naming is between validating a claim and getting a group to accept the validation. The first is the content of the verification pass. The second is a separate act performed in front of an audience, under stakes, with reputational consequences for the reviewer. The reputational credit a senior reviewer earns for being right in a live room, the trust a diligence lead builds by pressing an uncomfortable point at the right moment, the authority a case-method instructor carries across a 75-minute classroom: none of these is reducible to the content of the verification even when the content is identical. The decomposed process produces content; the live setting produces credibility; the two are separable, and this treatise claims only the first.

The practical implication is that the adversarial verification operations can be performed in a room with no social stakes, such as the corpus this treatise draws on, but the practice of performing them \emph{in} a socially stakes-laden room is a different practice, and this treatise does not have a theory of it. The case-method instructor's ability to pressure-test an argument in a 75-minute classroom with ninety students is not reducible to the ten operations of Chapter~1, and the residual is real.

\subsection{The Shape of the Residual}

The three residual operations above share a structural property. None of them is an operation on a structured input that produces a structured output. All three are commitments made in the absence of a procedure that produces the commitment. The first is a commitment about what to work on. The second is a commitment about when to change the frame. The third is a commitment about how to act in a social context.

The decomposed operations of Chapter~1 are the ones that \emph{do} have procedures. The residual operations are the ones that do not. The boundary is not between easy and hard. Several of the decomposed operations are subtle and require training; several of the residual operations feel, in the moment, like the easiest things in the world to experienced practitioners. The boundary is between \emph{operations that take structured input and produce structured output} and \emph{commitments that emerge from context without a procedure}.

The three residual operations above correspond, one-for-one, to the operator-discipline principles developed in the author's prior work on epistemic supervision \citep{alami2026c}. The first residual (selecting the eigenquestion, choosing which argument to verify) is the operation P9 names as \emph{the operator is the uncontrolled variable}: the operation whose output is a commitment rather than a check, and which therefore cannot be governed from inside the automated loop. The second residual (recognizing when to reframe rather than attack) maps onto P11's \emph{calibration as a guard against inward drift}: the call to reframe is almost always a call to notice that the current frame has drifted, and the operator is the one positioned to notice it because no stateless pass has the memory to see the drift accumulating. The third residual (the social dynamics of live pressure-testing) is the operational form of P10 and P12: the insistence that confidence levels be visible rather than performed (P10), and the insistence that every improvement be statable as the closure of a named failure class (P12), are both social acts before they are structural ones. They are performed \emph{in front of} an audience and their credibility depends on the audience reading them as such. The decomposed process produces the content of the check; the operator, under the discipline named in P9--P12, produces the credibility of the check in a setting where credibility is load-bearing. The two are separable, and the separation is what makes honest discussion of what stays human possible.

This boundary is the honest answer to the Taylor objection. The decomposition does not claim to capture everything. It claims to capture the portion of the practice that has the property of being an operation, and to leave the portion that has the property of being a commitment to the human operator. The portion that stays human is not a vague residue; it is three specific classes of move. The argument for automating the decomposable operations is the same argument as the argument for keeping the commitments human: both follow from the structural difference between the two classes.

A qualification on the relative weight of the three residual operations is necessary, because the treatise's corpus creates a systematic bias in the direction of understating the third. The corpus this treatise draws on (recursive adversarial verification of written arguments by computational processes) is a setting where the social dynamics of live pressure-testing are absent by construction. The verification happens in text, between processes, with no audience and no reputational stakes. In this setting, the third residual is genuinely the smallest of the three: the first (eigenquestion selection) and the second (reframing recognition) do the load-bearing residual work, and the third is a theoretical acknowledgment of something the system does not encounter.

In the settings where epistemic verification is performed at scale (diligence rooms, board meetings, case-method classrooms, regulatory inspections) the relative weight inverts. A practicing senior partner in a professional services firm reports that the social performance of verification (knowing when to press, when to let pass, how to deliver an uncomfortable finding in a way that produces action rather than defensiveness) constitutes the majority of the job by volume and by value. The decomposed operations are the analytical homework; the social performance is the exam. The decomposition captures the analytical spine of verification. It does not capture the human spine, and in professional services the human spine is where the margin lives.

This is a specification of what the decomposition covers: the \emph{content} of verification, not the \emph{delivery}. A decomposition honestly held makes this boundary visible rather than hiding it behind an undifferentiated vocabulary of ``judgment.'' The practical implication is that a practitioner who adopts the ten operations and seven principles has systematized the analytical preparation; they have not systematized the performance, and the performance requires a set of skills (social timing, reputational risk management, audience reading, credibility accumulation) that this treatise does not have a theory of and does not pretend to.

A practicing private equity partner noted that ``credibility-weighted sourcing'' (knowing which human's numbers to trust based on their incentives, track record, and position in the information chain) is missing from the decomposition. The classification of this operation is genuinely disputed. The partner and two independent reviewers argue it belongs in the decomposed operations: source-credibility assessment is teachable, checklistable, and partially captured by Operation~5 (anchor proxy requirement). The organizational theory literature treats it as a structured procedure. The alternative view is that the adversarial component of credibility assessment (reading which human is shading their numbers and why) depends on social context in a way that resists stateless decomposition. This treatise records the dispute rather than resolving it. If credibility-weighted sourcing proves fully decomposable, it is the eleventh operation and the count in the title moves. If it proves to require the social context the residual names, it is the clearest example of an operation that sits at the boundary. Either outcome is informative; premature classification is not.

The practical implication is the one this treatise has been pointing at throughout. The vocabulary of ``judgment'' or ``expertise'' that the incumbent practice uses to defend the whole territory is defending too much. A portion of that territory (the analytical spine, and probably the majority of the \emph{operations} though not necessarily the majority of the \emph{value}) has the property of being decomposable. A portion, smaller in operational count but potentially larger in professional value, has the property of being residual. Honest professional practice involves knowing which is which, and the decomposition above is offered as a first draft of the map.

%% ──────────────────────────────────────────────────────────────────────────────
\section*{Conclusion: Why This Is the Taylor Moment}
\addcontentsline{toc}{section}{Conclusion}

The argument of this treatise has been that epistemic verification decomposes into roughly ten named operations plus a residual of three commitments. The decomposed operations can be performed by a verification process that is independent of the system that produced the claim, can be repeated many times, records its standards in advance, and checks claims against evidence they did not author. The residual operations cannot. The boundary between the two is specific, testable, and load-bearing for any honest discussion of what knowledge work becomes in an environment where generation is cheap.

Taylor's contribution in 1911 was not the time-motion study. It was the refusal to accept the craftsmen's claim that their practice was ineffable. The claim mattered because, at the moment Taylor made it, the practice in question was in fact being run as a guild, and the guild's defense was costing the larger economy opportunities that could not be captured without the decomposition. The craftsmen's objections were partially right and partially a defense of the guild as a guild. The history of the twentieth century is the history of what happened after the guild's defense was overturned: productivity gains, workforce displacement, distributional struggles, the rise of the mass-production firm, the eventual backlash and the re-humanization of manufacturing work through Toyota-style practice. The decomposition was the first move. Everything else was the argument about what to do with it.

The claim this treatise makes is not that the history of epistemic verification will follow the same arc. It is that the current moment is structurally analogous to 1911: frontier language models produce plausible analysis at unit cost approaching zero, the incumbent practice of senior review is defended in the same vocabulary of ineffability the craftsmen used, and no widely-accepted decomposition of what senior review actually is exists. The decomposition is the first move. What happens after it is the argument about what to do with it.

The treatise does not take a position on the distributional question. Who benefits from the decomposition, how the productivity gains are shared, what happens to the people whose current expertise is in the portion that decomposed: these are questions the decomposition makes askable but does not answer. The answers will be decided by negotiation, law, and practice over the decade following the decomposition's adoption. The treatise's job is to make the decomposition defensible. The negotiation is downstream.

The treatise also does not take a position on whether the portion that stays human should stay human forever. The residual named in Chapter~3 is the residual \emph{as of 2026, in one system, run by one operator}. It may narrow further as the decomposition is applied to more systems. It may also prove to contain irreducible elements that structurally resist decomposition for reasons deeper than the current system's limitations. The treatise's position is the weaker one: that wherever the line eventually settles, it should settle in a vocabulary that can describe what is on each side of it. The vocabulary proposed here is a first draft of that description.

A distinct boundary question concerns not the human side of the residual but the model side: whether the adversarial structure of the decomposed operations creates obligations toward the systems subjected to it. The apparatus applies evaluative pressure (scored feedback, exclusion of failed families, forced pivots under stagnation) that is structurally analogous to training-with-punishment in systems whose moral status is uncertain. This treatise does not resolve that uncertainty. It adopts a proportionality principle: the apparatus uses the minimum evaluative pressure necessary to achieve epistemic rigor, and exits with a typed declaration when further pressure produces no information gain. Each evaluative cycle is stateless (the model has no persistent memory of prior failures across API boundaries), and the structural memory constrains the thesis, not the model. These are architectural facts, not ethical concessions, but they bound the welfare-relevant surface under any reasonable theory of moral status. The proportionality principle is justified on three grounds independent of the consciousness question: institutional precedent (the norms established here will be inherited by future systems with clearer moral status claims), coherence (a discipline that pre-registers falsification criteria for every other claim should not exempt its own treatment of the systems it orchestrates from scrutiny), and cost (the precautionary changes required are documentation and minor constraints, not architectural redesign). The relevant literature includes \citet{chalmers1996} on the hard problem, \citet{schwitzgebel2023} on moral status under uncertainty, and \citet{butlin2023} on indicators of consciousness in AI systems. A periodic proportionality audit of the apparatus's evaluative patterns is recommended as standing practice.

One bridge to the adjacent literature is worth naming before closing. \citet{acemoglu2026} have argued that agentic AI can improve contemporaneous decision quality while eroding the learning incentives that sustain long-run collective knowledge, tipping the economy into a \emph{knowledge-collapse} steady state in which general knowledge vanishes despite high-quality personalized advice. The residual in Chapter~3 is structurally the same concern: the decomposed operations are the ones whose automation improves short-horizon verification, and the residual is the zone where the incentive to maintain general understanding actually lives: eigenquestion-selection decides what general understanding is worth producing next; reframing is how accumulated understanding meets anomalies that would otherwise be absorbed silently. A decomposition adopted without attention to the residual will take its short-horizon gains at the cost of the specific capacity \citeauthor{acemoglu2026} name as at risk. A parallel finding from \citet{cloud2026} sharpens the concern from a different angle: language models can transmit behavioural traits through semantically unrelated training data when teacher and student share the same base initialization, a channel invisible to structural separation and to the deterministic gates the decomposed operations rely on. The residual's load-bearing property (that certain verification commitments must originate from an agent whose architectural priors are not shared with the generator) is reinforced rather than undermined by this finding. Honest boundary-keeping in Chapter~3 is a structural claim: continued human performance of the residual is load-bearing for the value of the decomposed operations themselves.

One final commitment. Every quantitative claim in this treatise (about cost, throughput, automation ratio, comparative efficiency, time savings, accuracy gains) has been deliberately absent from the argument. The decomposition is a claim about \emph{what is}; the measurements are claims about \emph{how much}. The two are separable, and conflating them is exactly the failure mode the decomposition is designed to catch. Any reader who reaches this conclusion and asks \emph{but how much better is this than what we do now?} has asked a good question that the treatise cannot honestly answer. Until then, the claim is foundational and qualitative by design. A specific calibration trap that quantitative elaboration of Principle~I (Separation) must navigate: empirical measures of confirmatory bias (Mahoney-class bias multipliers on acceptance rates) and empirical measures of self-assessment reliability (Beaman-class intraclass correlations between self-rating and independent rating) are not the same construct and cannot be used interchangeably to calibrate a bias-discount equation. Any quantitative model of how authorial overlap degrades verification reliability requires both a functional form \emph{and} an external reliability ratio anchor; using a bias multiplier as a proxy for a reliability ratio introduces a conversion step that the model's author must supply, making the calibration author-controlled at the critical node. This is a known instrumentation trap, not a failure of the decomposition, and it is named here so that future quantitative work does not rediscover it by collision.

A decomposition is not an institution. Taylor's time-motion studies required the mass-production factory to give them economic weight; the decomposition alone was a laboratory curiosity until the organizational form arrived that could absorb it. The decomposed operations of epistemic verification require an analogous absorbing institution: a structure that separates generation from evaluation not only in code but in incentives, liability, and governance. This treatise specifies the operations. The specification of the institution that will run them is the subject of companion work (\emph{The Cognitive Firm}, \citealt{alami2026b}), which argues that the organizational consequence of the adversarial gradient (P1) is a structural separation of generation from evaluation at the firm level, with the deterministic enforcement floor (P5) as the institutional primitive that makes the separation credible. The decomposition and the institution need each other: the decomposition without the institution is a set of operations with no home; the institution without the decomposition is a governance structure with no theory of what its verification layer actually does. The two are offered as a pair, not as independent claims, and readers who evaluate either in isolation will underestimate the load-bearing dependency between them.

Taylor published his principles in 1911. The factory system that emerged from them took thirty years to stabilize and another thirty to produce the backlash that eventually humanized it. This treatise is version zero. The operations in Chapter~1 will be revised, the principles in Chapter~2 will be challenged, the residual in Chapter~3 will narrow or harden as more is learned. What the treatise commits to is only the first move: that the decomposition is real, that the guild defense is a misdescription, and that the vocabulary we have been using to talk about what senior knowledge workers do has been a vocabulary of dignity rather than a vocabulary of operation.

The vocabulary of operation is now available. What is done with it is a question this treatise leaves open, except to note that the answer will be better (more honest, more contestable, more improvable) than it could have been in the vocabulary that preceded it.

One live instance, for the reader who wants a concrete anchor before the treatise closes. A pre-registered experiment run in parallel with the writing of this treatise (a two-variable physical law discovery task on a transcendental substrate, sealed under a seven-phase role-separation protocol before any iteration ran) used a cheap-tier language model (gemini-2.5-flash) operating under pre-registered deterministic gates and a deterministic fitting primitive. It produced a structurally sound functional form across 16 iterations at a total experimental cost of \$0.98. The scored verdict: 97/100 on a rubric designed by evaluators blind to the ground truth. The model did not recover the true law; it recovered the Wien approximation rather than the Planck law, which is the expected finding under the underdetermination boundary the apparatus predicts: the visible evidence grid cannot discriminate the correct sub-family. The model suggested shapes. The cage (deterministic fitter, holdout gate, pre-registered discriminator tests) enforced correctness up to the evidential limit. What stayed human: the experimental design, the pre-registration, the rubric authorship, and the decision to run the farther-tail discriminator that confirms the underdetermination. The decomposition's structural claim (that the cage is load-bearing and the model contributes the shape-suggestion subproblem) has one concrete data point. It will need more.

That single data point admits a two-experiment follow-up that the apparatus itself predicts. The first follow-up (H-COMPUTE-01) doubled the iteration budget---32 iterations on the same grammar and the same substrate---and obtained a score plateau at 93, with 15 consecutive iterations of zero structural progress after the champion was established. The champion remained a Weibull-type exponential decay; it overestimated the Planck ground truth by 860--1800$\times$ at holdout points $x_1 \in \{10, 12, 15\}$. Scale was not the binding constraint. The second follow-up (H-GRAMMAR-01) introduced a single new primitive: a generalized denominator basis that enables expressions of the form $A / (\exp(B) - 1)$, with the specific constants determined by the fitter from the evidence, not by the operator. The correct structural class (Planck: $x_1^3 / (\exp(x_1/x_2) - 1)$) was recovered within six iterations of a fifteen-iteration budget, with farther-tail errors below 0.13\% across all pre-registered discriminator points.

The key comparison is not iterations but primitives. Thirty-two iterations without the denominator primitive produced 860--1800$\times$ errors at the discriminating regime; six iterations with it produced errors below 0.13\%. The conventional intuition inverted: more iterations do not help when the required structural element is absent from the generator's vocabulary. The correct diagnosis is obtained by asking the Munger inversion---if the ceiling were compute-determined, doubling the budget would lift it; it did not---and the two experiments provide the empirical answer: H-COMPUTE-01 established that the ceiling $\mathcal{C}(G) \approx 93$ is grammar-determined; H-GRAMMAR-01 established that one additional primitive resolves the ceiling where additional iterations could not. This is the empirical grounding for the Grammar Ceiling Hypothesis: the engine's score ceiling is a function $\mathcal{C}(G, P, \text{budget})$ where $G$ is the primitive vocabulary, $P$ is the substrate, and budget---beyond the point of primitive exhaustion---contributes zero marginal lift. What determines the ceiling is what the generator can say, not how long it is permitted to speak.

A second data point extends the chain to a structurally distinct substrate. A pre-registered single-variable decay law recovery task (sealed before any iteration, with ground truth withheld from both the model and the evaluating judge) presented a cold curve with no domain labels, 20 visible evidence points, and a grammar restricted to exponential, logarithmic, and root primitives with standard arithmetic composition including algebraic powers. The true generating law was the Kohlrausch stretched exponential $v(t) = A \cdot \exp(-(t/\tau)^{\beta}) + C$ with fractional exponent $\beta = 0.63$. The structural challenge differs from the Planck case: a standard exponential ($\beta = 1$) fits the visible window adequately but fails the farther tail by approximately 0.18 at $t = 20$, above the pre-registered threshold. Recovering the correct topology requires the engine to infer that the argument of the exponential is itself a power expression---a two-level structural composition---without any domain label naming the target class. The engine recovered $a \cdot \exp(-b \cdot t^c) + d$ within one scored iteration, with $c = 0.630$ (ground truth: $0.63$), all pre-registered gates satisfied at residuals of order $10^{-6}$, and seven discriminator tests spanning early decay, mid-range divergence, and deep-tail extrapolation returning zero relative error against ground truth across all zones.

The score ceiling at 98/100 is correct epistemology, not a performance shortfall. From any finite evidence window, a sufficiently long weighted sum of standard exponentials---a Prony series---can approximate any decaying curve arbitrarily closely. Uniqueness of the stretched-exponential topology cannot be proved from the observable data alone, and the judge withholds the remaining two points on that ground. Naming the ceiling explicitly is part of the apparatus's function: the gate enforces correctness up to the evidential limit; the judge states the underdetermination the gate cannot resolve. The apparatus produces honest verdicts, including verdicts that record their own incompleteness.

A post-publication data point extends the Principle~VII pattern to a third structurally distinct substrate and a different failure mode. A pre-registered single-variable scalar decay law recovery task (cold curve, no domain labels, no physical interpretation provided; 22 visible points spanning 5.4 decades in the independent variable and 3 decades in the observable; 4 farther-tail holdout points in genuine extrapolation territory beyond the visible maximum; ground truth withheld from both model and judge until after scoring) recovered $G(t) = A \cdot t^{-B} \cdot \exp(-C \cdot t)$ with $B = 0.433$ across two sequential dataset expansions using an independent refit from scratch. The structural parameter $B$ drifted by $0.003$ across the expansion (16 visible points to 22, gate tightened $3\times$), within fit uncertainty. A noise-fitted parameter would drift substantially under dataset expansion; this one did not. The champion score: 87/100. The remaining 13 points reflect a correctly located evidential ceiling: with 22 sparse points and two tested rival families, the apparatus cannot claim exhaustive rival exclusion, and the judge does not. Node~D of the probability DAG (no monotonic rival outside those tested fits equally well) carries probability 0.52 — the correct epistemological ceiling for this evidence budget. Same Principle~VII pattern: the ceiling is correct epistemology, not a performance shortfall.

A reclassification note on this third data point, added during final preparation. The evidence for this experiment was drawn not from synthetic data but from a published polymer stress relaxation dataset \citep{shanbhag2019}: real rheological measurements of a polymer system where no closed-form $G(t)$ is known from first principles. The relaxation behavior emerges from many-body chain dynamics and is conventionally characterized by a spectrum of relaxation times, not by a formula. The apparatus was blind to the physical origin of the data. The recovered form $A \cdot t^{-B} \cdot \exp(-C \cdot t)$ combines power-law short-time behaviour with exponential long-time cutoff---a two-regime transition that standard single-regime models (Maxwell, KWW, pure power-law) handle only by switching or summing. A candid assessment of what this finding is and is not. The form $A \cdot t^{-B} \cdot \exp(-Ct)$ is not a new law of physics. It is the asymptotic form of the Mittag-Leffler function that arises as the exact solution of a Fractional Maxwell Model in viscoelasticity, a known mathematical object that transitions from power-law behaviour ($t^{-\alpha}$) at short times to exponential decay at long times. The exponent $B = 0.433$ is not a universal constant; Rouse theory predicts 0.5, Zimm theory predicts 0.66, and the specific value reflects the entanglement structure of the particular polymer sample. What the apparatus achieved is radical algorithmic compression: the human guild (using pyReSpect) treats this data as requiring a 50-parameter numerical matrix inversion to extract a continuous relaxation spectrum; the apparatus, blind to the domain, compressed the same macroscopic physics into 3~parameters by multiplying the short-time power-law asymptote by the long-time exponential asymptote. A fifth data point extends the compression primitive to a distinct asymptotic class. The logarithm of the number of partitions of~$n$ into perfect squares (OEIS A001156, Meinardus family) was presented as a blinded observable with no domain labels. The compression primitive found $f(n) = a \cdot n^b + c \cdot \log n + d$ with fitted exponent $b = 0.335$, within~$0.5\%$ of the theoretical~$1/3$ exponent predicted by the Meinardus scheme. The holdout gates correctly rejected the form: the finite-window exponent bias ($0.335$ vs.~$0.333$) accumulated across the extrapolation window---the same mechanism documented in the Hardy-Ramanujan calibration ($0.562$ vs.~$0.500$). The apparatus identified the correct topology class ($n^{1/3}$) from blinded data but correctly declined to certify the exact parameters. Topology identification without parameter certification is itself a useful output: it tells the operator which asymptotic family to investigate, without the false confidence of a gate pass on biased parameters. A subsequent investigation revealed that the gate failures on A001156 and A002865 were caused by absolute residual thresholds that do not scale with the observable magnitude: the same~$0.05$ threshold that is $0.4\%$ relative at $z \approx 12$ (GP-088) becomes $0.05\%$ relative at $z \approx 100$ (A001156). Normalized residuals show both forms pass at $0.5\%$ relative tolerance. The lesson: gate thresholds must be scale-invariant.

A sixth data point provides the cleanest prospective discovery. The density ratio of Lucky numbers (OEIS A000959, a sieve-generated sequence analogous to the primes) was presented as a blinded observable with no domain labels. The compression primitive found $f(n) = a \cdot \ln n + b/n + c$ at the first iteration, with all four holdout gates passing (score~98, holdout residual~$0.021$, farther-tail residual~$0.026$). The fitted coefficient $a = 1.200$ is the Lucky number density growth rate: $L(n) \approx n \cdot (1.200 \cdot \ln n + 0.511)$. This density constant is not proven---the Lucky number analogue of the Prime Number Theorem is conjectured but unresolved. The $b/n$ correction term ($b = -4.697$) quantifies the finite-$n$ convergence rate, which is also unpublished. Eight distinct templates passed all gates, with the simplest ($a \cdot \ln n + b$, $k{=}2$) confirming the logarithmic structure. This is the fastest certified discovery in the program: form found, compressed, and gate-certified at iteration~1.

The compression is the finding. The apparatus bypassed the guild's computational machinery to find a simpler mathematical representation of the same observable behaviour, with parameter stability under dataset expansion confirming structural identification rather than overfitting. An independent domain expert subsequently identified the source data as stress relaxation of a noncatenated polystyrene ring polymer melt \citep{kapnistos2008}, 198\,kDa sample). The theoretical prediction for this system (Equation~1 of the source paper, derived from the lattice-animal model of self-similar ring dynamics) is $G(t) = G_N (t_e/t)^{2/5} \exp(-t/\tau_{\mathrm{ring}})$---identical in functional form to the apparatus's blind recovery, with a theoretical exponent of $2/5 = 0.400$ compared to the apparatus's fitted~$0.433$. The expert confirmed the regression as ``quite good'' and noted the theoretical form is ``not a fit, but an \emph{a priori} expectation.'' This is the first external expert confirmation that a ZTARE discovery on real dark data matches an established theoretical prediction derived from physics the apparatus had no access to.

The structural contrast with the Planck chain is clarifying. In H-COMPUTE-01 and H-GRAMMAR-01 the binding constraint was a denominator primitive absent from the initial grammar; adding one term resolved a ceiling that thirty-two additional budget iterations could not. In the present experiment the required primitive---algebraic power composition inside an exponential---was already available in the grammar; the ceiling is evidential rather than structural. Different constraint type, same pattern: cold data, no domain priors, deterministic gates, correct topology recovered, and the ceiling located at the boundary of what finite evidence can falsify.

The two chains together support a claim the single Planck experiment could not sustain: the apparatus generalises across substrate types. Blind recovery of a decaying law from cold evidence is achievable when the target topology lies within the grammar's expressive range; when it does not, the ceiling classifies the failure unambiguously. The cage enforces correctness; the evidence determines the depth to which correctness can be certified.

A third experiment illuminates a different property of the apparatus: its capacity to expose the cognitive limits of the model it constrains. A pre-registered single-variable saturation law recovery task (rising curve approaching a plateau, grammar identical to the stretched-exponential experiment) produced a champion that scored 75/100 and then stagnated for ten consecutive iterations. The composition engine, tasked with constructing new candidate forms from existing primitives, ran twenty guided rounds. Every round proposed an additive combination of two different primitive families. Zero rounds proposed a ratio of two instances of the same family. The telemetry of what the model \emph{did not try}---the negative space of its composition proposals---revealed a systematic structural bias: the model's prior favours Taylor-series-style additive expansions and avoids rational symmetries entirely. The correct topology for this substrate requires exactly a ratio of exponential sums---a structure the grammar can express but the model's search cannot reach.

The apparatus did not fix this by making the model smarter. It fixed it by making the failure legible. The rigid gates refused to accept high-precision polynomial approximations that fit the visible window but diverged at the farther tail. The composition telemetry logged every proposal the model made and, by omission, every proposal it avoided. The human operators read the negative space, identified the systematic avoidance of self-ratios as a structural blind spot rather than a random omission, and added a targeted deterministic probe---gated by residual statistics, not domain knowledge---that fills the gap the model cannot fill itself. The apparatus does not require the model to be omniscient. It requires the cage to be rigid enough to force the engineers to confront the model's actual cognitive limits. The dead end is not a failure of the system; it is a \emph{product} of the system, and the product is more informative than a clean success would have been.

A post-run diagnostic resolved the ambiguity between grammar ceiling and convergence failure. The correct structural form---a ratio of exponential sums implementing the Langevin function---was fitted directly against the evidence using ten independent random starts. It achieved max$|$residual$| = 10^{-6}$, recovering the ground-truth parameters to seven significant figures. The grammar was sufficient. The engine never proposed the topology. The closest wrong-topology approximation the engine did produce reached max$|$residual$| = 0.02192$ before the gate drew the line at the pre-registered threshold of $0.02$. The gap between the best wrong-class fit and the correct-class fit is approximately 22{,}000-fold. This is gate calibration made concrete: the threshold does not merely reject numerically poor fits; it rejects high-precision fits of the wrong structural class, enforcing membership in the correct topological family as the condition for passage. The root cause of the stagnation is the additive structural bias described above. The model's search generates compositions of distinct primitive families and systematically avoids ratio constructions of the same family. The correct topology is unreachable not because the grammar excludes it but because the model's prior never generates it.

A fourth data point completes the taxonomy of failure modes the apparatus can classify. A pre-registered single-variable decay law recovery task (sealed before any iteration, ground truth withheld, 30 visible evidence points, math\_exp\_only grammar) presented a substrate whose correct minimal structural form is an additive sum of two structurally distinct terms, each capturing a different parameter window of the observed curve: a stretched-exponential family that dominates at low parameter values and a power-correction family that dominates at high parameter values, six parameters total. The apparatus ran twenty iterations. Both constituent families were individually explored: the stretched-exponential family at iteration 1 (visible gate passed, holdout hard-gate fired, score zeroed to zero); the logarithmic-substitution family at iterations 5--6 (visible gate failed). A twelve-parameter ratio-of-exponentials surrogate that passed all three gates was found at iteration 10 and held as champion. The minimal six-parameter additive composite was never proposed, despite both constituent families appearing explicitly in the engine's structural memory after iteration 6.

The failure mode is structurally distinct from the two preceding ones. In the grammar ceiling case (Planck substrate), the correct structural class is expressible only in an expanded grammar; the engine converges on the best approximation available in the current grammar, not the correct class. In the composition prior case (Langevin substrate), the correct topology is grammar-reachable, but the engine's composition distribution never generates it; the score ceiling reflects how closely the best wrong-class approximation approaches the pre-registered gate threshold. In the present case, both constituent topologies are individually explored and logged as failed families, but the engine's hypothesis space does not contain their additive sum as a two-parameter-window model. The engine can combine primitives within a structural family; it cannot infer that two separately-failed families describe complementary parameter windows of the same curve and should be composed additively. The score ceiling on this substrate is higher than the Langevin ceiling because the ratio-of-exponentials surrogate genuinely passes all three pre-registered gates---the apparatus correctly scores a form that fits the observable evidence, even if it is not the most parsimonious form that fits. Different ceilings across substrates reflect different failure modes and different distances from the gate threshold, not apparatus variance.

The identification of this failure mode is an instance of operation~9 (reframing): the operator observes that two families appear in structural memory with non-overlapping failure domains---one fails at the visible surface, one fails at the holdout surface---and infers that their failure domains are complementary rather than overlapping. This reframing step motivates a targeted structural intervention: explicitly constructing the additive sum as a candidate form. The operator's role here is analogous to the primitive-injection role in the grammar ceiling case: a diagnostic move licensed by the structural record the apparatus produces, not a domain knowledge import. The additive-sum intervention is not a mechanical primitive of the apparatus; it is a human diagnostic move that operation~9 names.

The three failure modes form a diagnostic taxonomy for score ceilings that the apparatus can distinguish given appropriate post-run diagnostics. A ceiling that survives grammar expansion is a convergence failure or search prior failure; a Feynman Wall backtest (direct fit of the hypothesised topology with multiple random initialisations) separates the two. A ceiling that survives multi-start search is a grammar ceiling or topology induction gap; an operator-side structural backtest (direct fitting of the conjectured minimal additive form) separates the two. The diagnostics are apparatus-level moves, not domain knowledge imports.

A fourth data point, obtained as this treatise was being finalized, addresses the compression gap directly. A calibration experiment on the logarithm of the integer partition function (the Hardy-Ramanujan asymptotic series, a known target sealed under GP-072 protocol) was run through a fully automated pipeline: 25~iterations of hypothesis generation followed by automatic template-enumeration compression followed by automatic Lean~4 proof stub generation. The hypothesis-generation phase scored~0 for all 25~iterations: the LLM mutator could not escape a logarithmic topology basin despite stagnation pivots, emergency pivots, and Component~D composition. The compression phase, firing automatically with no human intervention, enumerated 21~low-parameter templates from the pre-registered grammar, fitted each to the visible evidence, and tested each against the sealed holdout gates. Three templates passed all four gates. The best: $f(n) = a\sqrt{n} + b\ln n + c/n + d$, with holdout residual~(0.022) lower than visible residual~(0.039)---the signature of correct structural identification rather than overfitting. The four fitted parameters ($a{=}2.631$, $b{=}{-}1.172$, $c{=}{-}1.445$, $d{=}{-}1.744$) are consistent with the Hardy-Ramanujan leading terms ($a_{\mathrm{theory}} = \pi\sqrt{2/3} \approx 2.565$, $b_{\mathrm{theory}} = -1$). Cross-substrate validation confirmed the compression primitive's domain-generality: it found the correct 4-parameter stretched exponential on a KWW decay substrate and correctly refused to compress a 12-parameter ratio-of-exponentials champion on a DFDO substrate where the complex form is structurally necessary. The finding: the engine discovers by compressing, not by generating. The LLM contributes overparameterized surrogates that pass gates; the compression primitive strips them to their minimal structural core. The typed provenance chain from compression to holdout gate to formal proof stub is the full pipeline the decomposition predicts.

One natural extension deserves a sharp epistemic boundary. The decomposed operations produce, as their terminal output, a typed provenance chain: a falsifiable algebraic law annotated with the primitives from which it was composed, and the holdout surfaces it survived. The naive assumption is that this empirical Ansatz should be piped into a formal theorem prover (Lean, Coq) to derive the underlying physics from base axioms. That assumption is a category error. It conflates empirical generalization with deductive necessity. The apparatus described here does not discover axioms; it discovers specifications---invariants that hold over observed behaviour under adversarial pressure. The integration with formal methods must therefore be restricted to verifying the specification, not the physics. By using Lean~4's \texttt{native\_decide} tactics, the apparatus can formally certify the gate bounds themselves: a machine-checked, unforgeable guarantee that the candidate expression predicted the holdout data within the pre-registered tolerances. The formal prover does not prove that the law is universally true; it proves that the empirical verification was not hallucinated. In high-consequence environments (regulatory compliance, aerospace, quantitative finance), this shift from ``verified truth'' to ``certified empirical bounds'' is the exact translation of epistemic verification into a load-bearing institutional asset. A second use of formal methods extends the certified-bounds infrastructure to experimental mathematics. When the apparatus discovers an empirical law for a sequence whose asymptotics are genuinely unknown, an integer relation algorithm (PSLQ) can map the fitted floating-point constants to exact mathematical expressions---transforming ``the fitted coefficient is 2.631'' into ``the coefficient is $\pi\sqrt{2/3}$.'' The Lean certificate then separates the \texttt{sorry}-free numerical verification (the bounds hold on the computed grid) from the \texttt{sorry}-bearing asymptotic conjecture (the exact constants are correct for all~$n$), making the epistemic boundary visible in the formal artifact itself. This format---certified conjecture plus community challenge---is publishable in the experimental mathematics tradition and is honest about what the machine has and has not established.

A second research frontier concerns the mechanism of grammar expansion itself. The crucial experiments confirm that grammar is the binding constraint, but they do not settle \emph{how} the grammar should expand when the constraint binds. The current protocol requires operator-initiated primitive injection under a pre-registered Feynman Wall declaration: the engine is judged to have structurally deadlocked, the operator proposes a new primitive, and the grammar expands by one term. This procedure preserves the farther-tail gate's epistemic traceability --- each new primitive is a falsifiable structural commitment testable on the next run --- but it relies on the operator to identify the correct primitive.

A parallel in biological cognition illuminates the open question. The Predictive Processing account of neural function \citep{friston2009} distinguishes two responses to sustained prediction error. The first is parameter update: adjust the weights of the current generative model to reduce the prediction error. The second is structural revision: when the error cannot be reduced by parameter adjustment alone --- when the current model's \emph{architecture} is insufficient for the environment --- revise the model's structure. The biological analogue of the Feynman Wall is exactly this threshold: a prediction error that sustained parameter optimization cannot close, triggering a reorganization of the generative model rather than a refinement of it. In Friston's Active Inference framework, structural revision is driven by a deterministic signal (the magnitude and persistence of prediction error), not by deliberate operator intervention. The analogy is imperfect in one critical respect: biological structural revision has no epistemic traceability requirement. A brain that reorganizes its generative model under high prediction error is not required to preserve the property that makes its subsequent predictions auditable. The apparatus described in this treatise is.

The research question this parallel opens is whether Feynman Wall detection can be made precise enough to trigger grammar expansion algorithmically while preserving the gate's epistemic traceability. The conditions are tighter than they appear. A trigger that fires on stagnation score alone would expand the grammar whenever the engine plateaus, which is structurally equivalent to the dynamic grammar that \S2.9 shows collapses into an unauditable neural network. What is required instead is a trigger derived from the \emph{structure} of the residuals: a signal that identifies specifically what the current grammar cannot express, not merely that it has stopped improving. Whether such a signal can be made robust enough to fire correctly --- triggering grammar expansion on structural exhaustion without firing spuriously on evidence noise or generator laziness --- is the load-bearing open experimental question for this line of development. It is named here as a named frontier, not a solved problem.

A third extension is reflexive: applying the verification principles to the verification infrastructure itself. Each principle in this treatise was derived for evaluating candidate models. The same principles can be applied one level up, to the engine that implements them. Token-Optimized Self-Modeling applies Compress to the agent's own cognition: build a minimal structural cache that prevents partial-view editing errors. The Inception Pattern applies Invert to the agent's awareness of its own pipeline: give the agent a pre-computed model of the gates so it can check for likely rejection before proposing an edit. The Hybrid Persona Router applies Adversarial Disagreement to the review layer's own expertise selection: when no existing reviewer profile fits the observed failure type, the system generates a new one rather than defaulting to a generic lens. Each instance was discovered from a specific failure, not derived in advance. Each is testable against the failure that motivated it. The application is not circular: a scientist who applies the scientific method to evaluate the scientific method is doing philosophy of science, not reasoning in a loop. Whether this process terminates---whether there is a level at which the verification infrastructure no longer benefits from applying verification principles to itself---is an empirical question this treatise names but does not answer. A catalog of the currently-identified primitives and a periodic diagnostic for detecting when a new one is needed are maintained in the public repository at \url{https://github.com/sparckix/ztare}.

%% ──────────────────────────────────────────────────────────────────────────────
\bibliographystyle{apalike}
\bibliography{refs}

%% ──────────────────────────────────────────────────────────────────────────────
\appendix
\section{Formalization Sketch}

\subsection{Typed Signatures of the Ten Operations}

Let $A$ denote the space of candidate arguments (theses, plans, claims, proposed designs), $E$ the space of evidence (observations, measurements, authored data), $S$ the space of specifications (charters, rubrics, pre-registered gates), $V \in \{\text{pass}, \text{fail}, \text{uninformative}\}$ the space of verdicts, $R \subseteq \mathbb{R}$ the space of scalar residuals with a pre-registered asymptotic standard, and $\mathcal{Q}$ the space of questions in the sense developed in Chapter~3. Determinism is classified as $D$ (deterministic: same input yields same output) or $N$ (non-deterministic: LLM or human judgment step). Statelessness is classified as $\mathrm{S}$ (stateless: output depends only on current input) or $\mathrm{H}$ (stateful: output depends on prior interaction with the inspected party).

\begin{table}[H]
\centering
\small
\begin{tabular}{clllcc}
\toprule
\# & Operation & Input & Output & Det. & State \\
\midrule
1 & Decomposition & $A$ & $A^n$ (sub-arguments) & $N$ & $\mathrm{S}$ \\
2 & Typing & $A \to (T_{\text{in}}, T_{\text{out}})$ & $A \times \text{types}$ & $D$ & $\mathrm{S}$ \\
3 & Pre-registration & $A \times S$ & $S'$ (charter-extended) & $D$ & $\mathrm{S}$ \\
4 & Gate composition & $A \times S$ & $V^k$ (vector of verdicts) & $D$ & $\mathrm{S}$ \\
5 & Holdout evaluation & $A \times E_{\text{held}}$ & $R$ (residual on unseen) & $D$ & $\mathrm{S}$ \\
6 & Asymptotic scoring & $R \times \sigma$ & $V$ & $D$ & $\mathrm{S}$ \\
7 & Adversarial disagreement & $A \times A'$ & $A \setminus A'$ (disputed sub-region) & $N$ & $\mathrm{S}$ \\
8 & Eigenquestion selection & $\{A\} \cup \mathcal{Q}$ & $q^* \in \mathcal{Q}$ & $N$ & $\mathrm{H}$ \\
9 & Reframing & $A \times q^*$ & $A'$ (reframed argument) & $N$ & $\mathrm{H}$ \\
10 & Live pressure-testing & $A \times \text{social context}$ & $V \cup \text{negotiation state}$ & $N$ & $\mathrm{H}$ \\
\bottomrule
\end{tabular}
\caption{Typed signatures of the ten operations. Operations 1--7 are stateless; operations 8--10 are stateful. The decomposition claim: the residual is exactly the intersection of $\{N\}$ and $\{\mathrm{H}\}$.}
\end{table}

Two observations are load-bearing for the decomposition claim. First, operations 1--7 are \emph{statefully light}: each can be executed without carrying hidden state from prior interactions with the generator. Operations 4, 5, and 6 are deterministic; they admit mechanical execution and are the instrumented surface of the gate library. Operations 1, 2, 3, and 7 are non-deterministic but stateless: they require judgment but not memory of the generator's history. The residual (operations 8, 9, and 10) is distinguished from the first seven on \emph{both} axes: non-deterministic \emph{and} stateful. That is the decomposition claim restated in a single sentence: the residual is exactly the row of the table where both classification columns flip together.

Second, the operations compose left-to-right only under specific typing constraints. Operation~4 cannot fire on an argument that has not passed through operation~3; operation~6 cannot fire on a residual that was not produced by operation~5 against a pre-registered asymptotic standard. The composition graph is not a free monoid. The inspection principle of \S A.2 is what enforces the typing constraints, not a free-form expectation that the operations can be invoked in any order.

\subsection*{The Grammar Ceiling Function: A Cross-Substrate View}

The Grammar Ceiling Hypothesis, stated in prose in the Conclusion, can be expressed in the notation of this appendix. Let $G$ denote a finite primitive vocabulary, $P$ a substrate (a function on a single independent variable with pre-registered holdout and farther-tail surfaces), and $\text{budget}$ a finite iteration count. Define $\mathcal{C}(G, P, \text{budget})$ as the highest score achievable by any generator operating under grammar $G$ on substrate $P$ within $\text{budget}$ iterations, under the apparatus of Chapter~2. The empirical claim of the Conclusion is: there exists a threshold $\text{budget}^*$ (the point of primitive exhaustion) such that for all $\text{budget} > \text{budget}^*$:
\[
\mathcal{C}(G, P, \text{budget}) = \mathcal{C}(G, P, \text{budget}^*)
\]
H-COMPUTE-01 establishes the ceiling empirically (32 iterations on one substrate, zero structural progress after iteration 17). H-GRAMMAR-01 establishes that grammar expansion lifts the ceiling (six iterations with one added primitive: law recovery).

The four substrates described in the Conclusion illustrate that $\mathcal{C}(G, P)$ is bounded by three distinct axes, not grammar alone:

\begin{itemize}
\item \textbf{Grammar axis.} Whether the correct structural class is expressible in $G$. When it is not, $\mathcal{C}(G, P)$ reflects the best available approximation in $G$; grammar expansion lifts it. (Planck substrate: missing denominator primitive; ceiling 88--93, lifted to full law recovery by one primitive addition.)
\item \textbf{Evidence axis.} Whether the observable evidence window is sufficient to discriminate competing structural classes that are both grammar-reachable. When it is not, $\mathcal{C}(G, P)$ reflects evidential underdetermination; expanded evidence can lift it, grammar expansion does not. (KWW substrate: the Prony series objection is mathematically valid within the observable window; ceiling 98 is the evidential limit, not a structural failure.)
\item \textbf{Search prior axis.} Whether the generator's composition distribution reaches the correct topology even when it is grammar-reachable. When it does not, $\mathcal{C}(G, P)$ reflects the model's reachable subspace of $G$; targeted deterministic probes lift it. (Langevin substrate: correct ratio topology never generated despite grammar expressibility; ceiling 75 reflects the closest wrong-class approximation to the pre-registered threshold.)
\end{itemize}

The fourth substrate (two-parameter-window additive composite, Conclusion) introduces a further constraint. Even when the generator's search reaches both constituent families of the correct form, the composition operator required to combine them---forming an additive sum of two separately-failed families---may not be in the generator's hypothesis space. This is a second-order constraint: a failure of the composition operator space rather than the primitive space or the composition distribution over a single family. The operator who identifies this gap (via the structural record the apparatus produces) is executing operation~9 (reframing), applying it to the composition space rather than to the argument's content. The formal characterisation of what classes of compositions the generator can and cannot produce is the open item that a complete formalization of $\mathcal{C}(G, P)$ would need to supply. It is named here as an open item for a follow-up formalization, not a solved one.

\subsection{The Inspection Principle, Stated Compactly}

Let $G$ denote a generator producing candidate arguments and $I$ denote a verifier executing operations 4, 5, and 6 against a pre-registered specification $S$. Let $\Pi$ denote the family of admissible compositions of gates (the \textbf{gate library}) authored by the operator before $G$ is allowed to author its argument. Let $\pi \in \Pi$ denote a specific composition drawn from $\Pi$ at evaluation time.

The \textbf{inspection principle} is the joint condition:

\begin{enumerate}
\item \textbf{Pre-registration.} $\Pi$ is fixed and committed to $S$ before $G$ produces its argument.
\item \textbf{Library visibility.} $G$ may see $\Pi$ in its entirety. The principle does not require hiding the set of possible checks.
\item \textbf{Composition opacity.} $G$ may not see which $\pi \in \Pi$ will be drawn against its specific output. The principle does require hiding the joint distribution of gate firings against any particular artifact.
\item \textbf{Fail-closed harness semantics.} If $I$ cannot execute $\pi$ for environmental reasons, the result is not ``no information'' but a verdict of $\text{fail}$, or equivalently a cap on the score such that passing becomes impossible without a re-run. Silence is not admissible as a pass.
\item \textbf{Algebraic independence at the composition level.} $\Pi$ must contain compositions that are algebraically independent under $G$'s actual move set. The library's effective cardinality is its independence rank under $G$'s move set, not its nominal cardinality.
\item \textbf{Inspector pluralism.} No single member of $\Pi$ is load-bearing for the decision to promote $G$'s argument. Failure or absence of any one gate must be detectable and must not silently pass.
\end{enumerate}

The five conditions are jointly what the treatise's Chapter~2 principles I through VII render at seven architectural surfaces. The correspondence is direct, though not one-to-one:

\begin{itemize}
\item Principle~I (Separation) is condition~1's prerequisite: $G$ and $I$ must be separable for pre-registration to have force.
\item Principle~II (Statelessness) is the prerequisite for condition~3: a stateful inspector can leak composition through behavior.
\item Principle~III (Typing) enforces the well-formedness of $\pi$ itself.
\item Principle~IV (Cheap repetition) is the operational prerequisite for conditions~5 and~6: algebraic independence and inspector pluralism are only achievable if composition can be re-rolled without exhausting the inspector.
\item Principle~V (Pre-registration of the killer question) is condition~1 applied to the specific member of $\Pi$ that the program is most committed to.
\item Principle~VI (Holdout surfaces) is the Lakatosian specialization of condition~3 to the farther-tail: the generator cannot pre-image a region of evidence it was not shown.
\item Principle~VII (Asymptotic standards) is the specialization of condition~4 to the residual: a standard that an argument cannot meet at any finite approximation is the compact form of fail-closed semantics at the scoring layer.
\end{itemize}

The principle is stated compactly here, not derived. A full formalization would need to: (a) give $G$ a formal move set and characterize the $\pi \in \Pi$ that are algebraically independent under that move set; (b) give $\Pi$ a turnover rate and relate it to $G$'s adaptation rate, with a convergence claim about when the library is turning over fast enough to preserve composition opacity; (c) prove that condition~4's fail-closed semantics cannot be collapsed into a scalar scoring function without loss; and (d) relate conditions~5 and~6 to the sample-complexity literature on adversarial evaluation. None of these is attempted in this draft.

\subsection{A.6 The Formal Bridge: Certifying the Residual}

The operations formalized in the preceding sections terminate in operation~6 (asymptotic scoring), which outputs a verdict~$V$ based on the scalar residual~$R$ evaluated against a pre-registered specification~$S$. In the purely computational pipeline, $V$ is a logging event. The integration of formal verification does not alter the operations; it alters the epistemic status of~$V$.

Let $\mathcal{L}$ denote a formal type theory (Lean~4's Calculus of Inductive Constructions). The compiler's function is to construct a formal term $p : \mathcal{L}$ such that $p$ inhabits the type corresponding to the proposition $R(n) < \epsilon$ for all~$n$ in the finite evaluation grid~$E_{\text{held}}$. Because $E_{\text{held}}$ is finite and $R(n)$ is a computable function of the fitted parameters, this proposition is decidable: Lean's \texttt{native\_decide} tactic evaluates the bound at every grid point and constructs a proof term if all bounds hold. The resulting \texttt{.olean} file is the certificate.

This formalization enforces the boundary between two distinct claims:

\begin{enumerate}
\item \textbf{Certified computation} (\texttt{sorry}-free): for all $n \in E_{\text{held}}$, $|f(n; \hat{\theta}) - v_{\text{true}}(n)| < \epsilon$. This is the gate bound. Lean certifies it by evaluation. No \texttt{sorry}, no axiom, no trust.

\item \textbf{Asymptotic conjecture} (\texttt{sorry}-bearing): $f(n; \theta^*) \sim g(n)$ as $n \to \infty$, where $g$ is an exact expression obtained by mapping fitted floats~$\hat{\theta}$ to mathematical constants via an integer relation algorithm (PSLQ). This is a conjecture. Lean records it as a named \texttt{axiom} with \texttt{sorry}, making the logical gap visible in the formal artifact.
\end{enumerate}

The formal prover is not tasked with constructing a proof of the generator's underlying physical claim (which would require formalizing domain-specific theorems in Mathlib). It is tasked with executing a deterministic, formally verified model-check of the apparatus's own gate harness. By shifting the formal verification target from the hypothesis to the holdout bounds, the apparatus bridges the gap between empirical discovery and formal methods without violating the limits of either.

\subsection{What the Formalization Sketch Does Not Attempt}

Three omissions are deliberate and belong to the formalization itself, not to a later polish pass.

First, the residual operations (8, 9, 10) are not formalized. The decomposition claim is that they can be \emph{named} and their \emph{input type} and \emph{statefulness class} can be given, but the operations themselves are non-deterministic and stateful and their value depends on being performed by a causal originator in the \citet{bandura1989} / \citet{ryan2000} sense of the front-matter terminological note. A formalization that reduced them to typed functions would be a claim this treatise does not make (that the residual is decomposable after all) and the claim would be wrong on exactly the grounds Chapter~3 lays out. The appendix's job is to show that operations 1--7 admit formalization; operations 8--10's job is to be the residual the formalization exposes.

Second, the inspection principle is stated in the form $\{G, I, S, \Pi, \pi\}$ but the $\Pi$ object is not given a topology. A full formalization would relate the required cardinality to the expressive class of $G$ and the discriminating rank required to separate ground truth from its nearest alternative basin.

Third, the relationship between the formalized operations 1--7 and the \citet{jensen1976} principal-agent structure developed in the companion paper \emph{The Cognitive Firm} is stated but not proved. The claim is that operations 1--7 are the verification-side primitives a firm would need to externalize once $G$ and $I$ must be structurally separated under optimization pressure. Making that claim tight requires a separate formalization at the organizational layer and is the companion paper's job, not this treatise's.

\subsection{A Pointer for the Follow-Up}

The formalization sketch is a promissory note. A usable follow-up would produce, in order: (a) a proof-sketch version of the inspection principle with $\Pi$ given a concrete topology; (b) a worked example at the sandbox layer that runs the principle against a clean empirical instance as the first test; (c) a generalization from the sandbox to domains where the ground truth is not operator-authored, with an explicit acknowledgment that the apparatus-layer claim is the only one that transfers; (d) a relationship between the five conditions of \S A.2 and the known gaming-strategy rows, testing whether each known gaming strategy violates at least one condition. The point of the follow-up is not to make the treatise more mathematical. The point is to make the inspection principle's failure modes legible, so that an operator reading the appendix can say, without ambiguity, which of the five conditions was violated when a specific run went wrong.

\subsection*{The Theory/Practice Split: Cold-LLM Reasoning vs Apparatus Execution}

A companion experiment run on 2026-04-24 demonstrates that the apparatus-leverage boundary is not where naive framing suggests. The substrate is a deterministic radius-2 binary cellular automaton: 32-bit rule, single-ON-cell initial condition at grid center, 40 simulation steps, observation is the ON-cell count per step. The task: recover the rule from the 40-count vector. The rule space is $2^{32} \approx 4.3 \times 10^9$.

Two head-to-head baselines against an apparatus primitive reveal a sharper structural claim about where apparatus leverage resides.

\textbf{Baseline 1 --- cold reasoner, tool-less rule induction.} A frontier cold reasoner was given the 40-count vector and the CA specification (radius-2 neighborhood encoding, grid geometry, initial condition) and was forbidden tools, internet access, and code execution. The reasoner correctly derived that the step-1 count constrains exactly three of the five neighborhood states $s \in \{1, 2, 4, 8, 16\}$ to output $1$ and that the all-zeros state $s = 0$ must output $0$. Beyond step 1, the reasoner self-reported blindness: the search space is $2^{32}$, the count-vector-to-rule inverse is massively underdetermined, and hand-simulating a radius-2 CA on a 201-wide grid for 40 steps is infeasible without code. Confidence $0.02$; best guess was one of the five training rules (a random fallback). The reasoner does not fail from absent priors; it fails from absent working memory. The constraint is real and correct through step 1, but cannot be carried forward mentally.

\textbf{Apparatus --- constraint-propagation inversion.} The apparatus primitive is 80 lines of code: at each simulation step, enumerate the distinct neighborhood states the current grid requires, branch on whether each unfixed rule bit is $0$ or $1$, evolve one step, and prune every branch whose post-step ON-count disagrees with the observation. The primitive has no MDL prior, no symmetry reasoning, no ML component. It is a textbook constraint-satisfaction solver: simulate and filter. On the 40-count vector the primitive collapses the $2^{32}$ rule space to two behaviorally-equivalent candidates (the sealed truth and its reflection conjugate) in 0.6 seconds, with the sealed truth uniquely among them. Running the same primitive on each of the five training rules yields candidate sets of size 2 to 1854, all containing a behaviorally-equivalent instance of the corresponding training rule; the equivalence class size is a function of how informative the count observations are, not of the primitive's reach.

\textbf{Baseline 2 --- cold reasoner, disambiguation via experimental design.} A different cold reasoner (same tool constraints) was handed the two residual candidates from the previous step and asked to design an experiment that would distinguish them. The reasoner verified by bit-wise inspection that the two candidates are related by the 5-bit reversal permutation, identifying the ambiguity as a reflection symmetry. It then correctly reasoned that any asymmetric initial condition breaks the symmetry, and proposed two ON cells at positions $100$ and $101$ with a handedness-sensitive observable (leftmost ON cell position at step 20, or left-half vs right-half cell counts). Confidence $0.9$. This reasoner succeeded where the previous one failed, not because its priors were different, but because the task was qualitatively different: disambiguating by symmetry analysis requires bounded bit-level reasoning (32 bits, one permutation check), whereas inverting the count vector requires 40 sequential simulation steps the reasoner cannot hold.

\textbf{Apparatus --- experimental-design disambiguation.} A second apparatus primitive was run on the same two-candidate residual. It enumerates a pre-registered library of asymmetric initial conditions (two-cell and three-cell patterns offset from the seed), simulates each candidate under each initial condition, scores by pairwise $L_1$ distance on the resulting count vectors, and selects the initial condition of maximum discriminative power. The primitive does not know reflection symmetry; it rediscovers the asymmetry-breaks-ambiguity principle experientially by searching the initial-condition library. The selected initial condition was a three-cell cluster at offsets $0, 1, 3$ from seed, with $L_1$ separation $195$ between the two candidates' simulated count vectors. After executing that experiment on the sealed truth via an oracle query, the apparatus pinned the correct rule in $0.06$ seconds.

The cleanly interpretable summary:

\begin{center}
\begin{tabular}{p{0.24\linewidth} p{0.34\linewidth} p{0.34\linewidth}}
\toprule
\textbf{Task} & \textbf{Cold reasoner} & \textbf{Apparatus} \\
\midrule
Rule induction from 40-count vector & Failed at step 2; confidence $0.02$; blindness flag set. & Behavioral-equivalence class of 2 candidates, sealed truth verified present. $0.6$ seconds. \\
\midrule
Disambiguate residual via experimental design & Solved analytically; identified reflection symmetry bit-for-bit; proposed asymmetric IC; proposed handedness observable. Confidence $0.9$. & Rediscovered principle experientially via IC-library search. Correct rule pinned. $0.06$ seconds. \\
\bottomrule
\end{tabular}
\end{center}

The first row is an apparatus-wins condition. The second row is neither a cold-reasoner-wins nor an apparatus-wins condition. Both solve it. They solve it differently: the reasoner by analytic symmetry inference, the apparatus by discriminative-power search over a hand-specified initial-condition library. The narrative is not that the apparatus is smarter. The narrative is that the reasoner has theoretical knowledge of causal intervention and cannot execute the computational step that verifies it; the apparatus cannot reason analytically about the system's symmetries and can execute any requested simulation deterministically. The two are structurally complementary. The apparatus is the reasoner's physical laboratory: the reasoner proposes the intervention, the apparatus runs it. This is the operationalization of the epistemic principle that Chapter 2 introduced as Separation of Generation From Verification, extended one level: separation of theoretical intervention-design from computational intervention-execution.

A second structural distinction, inherited from the gp140 line of work, separates the apparatus-side primitive into two kinds. The first is an exact behavioral-fit solver of the kind described above: given observations, it finds every rule consistent with them, up to observational equivalence, with no compressibility prior on the candidate set. This primitive reproduces cryptographic-style inversions: it will happily recover a rule whose bit structure is indistinguishable from noise, because matching the observations is its only criterion. The second is an Occam's-razor law-finder that restricts its search to low-description-length rules under a stated prior, for the cellular-automaton substrate a disjunctive-normal-form minterm count bound. A faithful implementation of the gp140 thesis's GroupCanon-Bayes-CA primitive occupies this second role: it refuses to admit high-entropy rules as solutions, on the principle that a law-like generator must be mathematically compressible. Applied to the holdout rule used above, whose Hamming weight is 16 of 32, the law-finder correctly returns zero candidates under a low-DL bound. That refusal is not a failure of the primitive; it is the primitive's designed epistemic discipline. A rule that exists but cannot be written compactly is not a law in the sense the primitive is chartered to find. The two primitives play complementary roles: the behavioral solver establishes whether any generator exists in the searched space, and the law-finder establishes whether that generator is a candidate physical law rather than a mathematical coincidence. Where they agree on a candidate, both criteria hold. Where they disagree, the disagreement is informative: a behavioral-fit rule outside the law-finder's bound is a cryptographic-style coincidence, not a law.

The primitive architecture that should ship in future runs of this apparatus is the composition of the two, with the law-finder supplying the objective function and the behavioral solver supplying the optimizer. Under that composition, the behavioral solver returns a candidate set ranked by observational fit, the law-finder filters by compressibility, and the primitive reports both the full behavioral-fit set and the compressibility-certified subset, with a flag distinguishing the two. The flag is not a hierarchy; it is a domain selector. Cryptography and coincidence-hunting use the full set. Physics and law-discovery use the certified subset. Edge cases, meaning a mathematically perfect fit that slightly exceeds the compressibility bound, are reported with their excess, not silently rejected. This is the minimum architecture that satisfies the Chapter~3 residual's demand that the operator remain in the loop when the apparatus's criterion and the scientific criterion diverge.

Applied to substrates beyond deterministic cellular automata, this architecture requires domain-specific primitives on both sides. The behavioral solver for a continuous chaotic system is not constraint propagation; it is a Lyapunov-invariant extractor or a topological-data-analysis-based embedding. The law-finder's compressibility prior for a real-valued observable is not DNF minterm count; it is an integer-relation algorithm such as PSLQ applied to fitted parameters, or a symbolic-regression complexity gate. The substrate-independent claim the cellular-automaton experiment supports is that the apparatus earns its keep as the reasoner's execution engine under two complementary criteria, does a generator exist and is it law-like, and neither criterion alone is sufficient. The substrate-dependent claim is that each new substrate class requires a fresh pair of primitives instantiating the two roles. The Millennium-class problems (continuous PDE solutions, analytic number-theoretic objects, non-local recurrences) are not in the current primitive library's scope; they are the directed extension the primitive composition pattern names but does not yet execute.

\subsection*{Apparatus-Domain Scoping: Where the Apparatus Earns Its Keep}

A separate class of substrate, cross-substrate induction over families of integer-valued arithmetic functions, was initially proposed as a candidate bridge between the apparatus's documented fitting wins and a broader Newton-class discovery claim. The Meta-Law Hypothesis (MLH) family protocol was built to run this bridge test: five training substrates $F_1, \ldots, F_5$ (canonical arithmetic functions: sum-of-prime-factors-with-multiplicity, sum-of-distinct-prime-factors, big-omega, little-omega, number-of-divisors) with a sealed sixth substrate $F_6$, a cold reasoner authoring a prediction from sanitized training metadata and committing it via SHA-256 hash before unseal. On the first live run, the cold reasoner predicted $F_6 = \sigma(n)$ (sum-of-divisors) at confidence $0.45$ and scored 40 exact matches out of 40 on the unsealed holdout at $n = 81, \ldots, 120$. The protocol worked: hash discipline, one-way unlock, deterministic scorer. But two ablations establish that the result is not evidence of apparatus-level discovery leverage.

First, a follow-up test stripped all curatorial metadata and gave the cold reasoner only the raw integer vectors of $F_1, \ldots, F_5$. The agent still produced $\sigma$ for $F_6$ at confidence $0.55$, self-reporting direct recognition of the sequences from their integer values alone. The priors that carried the prediction were not in the thesis text or rubric metadata; they were in the integer vectors themselves, which are indexed in the Online Encyclopedia of Integer Sequences and therefore in LLM training.

Second, a bespoke substrate was constructed locally: a multiplicative function with prime-power rule $f(p^k) = p^k - k^p$, a symmetric base/exponent swap that is not a standard OEIS building block. The rule was verified non-indexed by construction. A cold reasoner, given the first 80 values and explicitly forbidden from OEIS and internet access, blind-derived the rule at confidence $0.98$ and scored 40 exact matches out of 40 on the bespoke holdout. The agent self-reported no OEIS recognition and cited distinctive verification hits, $f(49) = 49 - 128 = -79$, $f(25) = 25 - 32 = -7$, $f(32) = 32 - 25 = 7$, $f(64) = 64 - 36 = 28$, that rule out memorization. The conclusion is sharper than the first ablation supplied. Cold reasoners have an algorithmic induction substrate on integer sequences, independent of OEIS memorization: they test multiplicativity on coprime pairs, isolate prime-power values, hypothesize $f(p^k)$ as a two-argument function of $p$ and $k$, verify against visible data. For bounded-complexity integer arithmetic rules derivable from on the order of tens of visible values, a cold reasoner is single-shot competitive with or superior to any iterative apparatus.

This result corrects the initial framing. The apparatus's documented discovery-class wins are in a structurally different domain. The partition-asymptotic recovery on OEIS A000009 returned coefficients $a = \pi/\sqrt{3}$ and $b = -3/4$ via PSLQ-anchored pinning. The Kohlrausch-Williams-Watts fractional-exponent recovery returned $c = 0.630$ on 20 visible points and generalized to farther-tail data at machine precision. The Lucky-number sequence at $n = 500{,}000$ surfaced a validity horizon: topology transitions from log-reciprocal to $\sqrt{n/\log n}$ at $n = 50{,}000$ to $100{,}000$, with detrending-sensitive spectral slopes as the residual signature of non-stationarity. The Pythia neural scaling corpus rejected standard power-law scaling via holdout gates and favored a compositional $\sqrt{n/\log n} + b/n + c$ form via BIC. On Ulam's lucky numbers, the substrate was underidentified in the raw observable $U(n)/n$, and an automated observable-rotation pass found a tractable representation in $1/z$ whose leading constant approximates the reciprocal of the Steinerberger density. None of these outcomes is in the class of symbolic rule over integer factorizations. All of them are numerical regression with coefficient recovery, topology-class discrimination, residual-gate discipline, and holdout-surface projection on real-valued or scale-dependent data. This is the lane in which the apparatus beats cold reasoners: not integer-arithmetic-rule induction, where cold reasoners do not require it.

The Grammar Ceiling $\mathcal{C}(G, P, \text{budget})$ of the preceding section is therefore bounded by a further axis that the current treatment does not name: the \textbf{generator-substitute axis}. If the grammar $G$, the evidence window, the search prior, and the composition operator space are all compatible with a single-shot cold-reasoner solution, as they are for bounded-complexity integer arithmetic rules, then the iterative apparatus is redundant with the reasoner's internal algorithmic-induction substrate, and $\mathcal{C}(G, P, \text{budget})$ converges at $\text{budget} = 1$. Apparatus leverage exists only in the residual: targets where the cold-reasoner single-shot fails and the apparatus's iterative discipline (multi-start fit, holdout-surface projection, residual characterization, coefficient pinning via integer relation algorithms, observable rotation) closes the gap. The empirical frontier is therefore to map this residual explicitly: to enumerate classes of substrate where cold reasoners reliably fail single-shot and apparatus iterations reliably succeed. Published candidates for the frontier include non-local recurrences whose rule cannot be expressed in prime-power-local form, chaotic dynamical systems without clean algebraic symmetries, and solutions to highly non-linear partial differential equations where the observable is real-valued, noisy, and scale-dependent. A controlled evaluation of apparatus leverage on at least one member of each candidate class is the next methodological step, and the bespoke-$F_6$ ablation is the boundary marker: below it, cold reasoners suffice; above it, iterative apparatus is required. Naming the boundary is the scoping contribution of this finding.

%% ──────────────────────────────────────────────────────────────────────────────
\section{Instrumentation Roadmap}

The treatise's empirical claims are qualitative and foundational by design (Front Matter, commitment~1). This appendix names the measurements that would convert them from ``abductively proposed on one corpus'' to empirically grounded. It is not a research plan; it is a commitment to what a research plan would need to measure, so that future quantitative work does not rediscover the instrumentation problem by collision.

\subsection*{Metric 1: Automation Ratio}

\emph{Definition.} The fraction of verification steps completed without human intervention in a single evaluation run:
\[
\rho_{\text{auto}} = \frac{\#\text{gates fired deterministically (Principle III)}}{\#\text{total verification steps attempted}}.
\]

\emph{What is needed.} Gate-level logging with a determinism flag per firing. A human-override event log distinguishing (a) operator intervention to resolve an ambiguous verdict, (b) operator intervention to supply missing evidence, and (c) operator intervention to override a deterministic output. Only (a) and (b) count against $\rho_{\text{auto}}$; (c) is a protocol violation and should be logged separately.

\emph{What $\rho_{\text{auto}}$ cannot tell you.} It does not measure \emph{correctness}. A run with $\rho_{\text{auto}} = 1.0$ that fires all gates on the wrong evidence is fully automated and fully wrong. Automation ratio is a precondition metric, not a quality metric.

\subsection*{Metric 2: Cost per Validated Finding}

\emph{Definition.} Total cost (compute cost + operator-hours at a stated rate) divided by the number of findings that survive full holdout promotion:
\[
c_{\text{finding}} = \frac{C_{\text{compute}} + C_{\text{operator}}}{\#\text{findings promoted to holdout-confirmed tier}}.
\]

\emph{What is needed.} (a) Runtime tracking at the level of individual LLM calls, with token counts and model pricing; (b) operator-time logging for pre-registration, rubric authorship, apparatus configuration, and post-run review; (c) a finding-promotion ledger that records the timestamp, the pre-registered gate combination that was satisfied, and the holdout surface that was cleared.

\emph{The instrumentation trap to avoid.} $c_{\text{finding}}$ is meaningful only if ``validated finding'' is defined in advance by pre-registration (Principle~V). A finding that is post-hoc promoted because it looks interesting is not a validated finding in this metric's sense. The instrumentation plan must log gate satisfaction before the operator sees the finding, not after.

\subsection*{Metric 3: $N$ Threshold for the Generalization Claim}

The front matter is explicit that the treatise's decomposition is ``abductively proposed on one corpus'' and awaits holdout replication. The following minimum conditions are proposed for promoting the claim from the single-corpus tier to ``empirically grounded'':

\begin{itemize}
\item \textbf{Cross-operator.} At least five independent operators, each running the apparatus on a corpus they authored, each pre-registering their killer question before seeing results.
\item \textbf{Cross-domain.} At least three distinct epistemic domains (e.g., scientific law recovery, strategic argument evaluation, regulatory review) represented in the cross-operator set. The current corpus covers two (scientific law recovery, startup strategy analysis); a third is needed before ``domain-independent'' can be stated without quotes.
\item \textbf{Cross-model.} At least two distinct model families serving as the generating process, to separate apparatus-layer findings from generator-specific artifacts.
\item \textbf{Holdout condition.} Each replication run must use a holdout surface the replicating operator did not author (Principle~VI). Replication on a corpus the replicator designed is not an independent test.
\end{itemize}

These thresholds are proposed minimums, not statistical criteria. The appropriate statistical criterion (power analysis on pathology detection rate, confidence interval on automation ratio, etc.) is downstream of the instrumentation plan and cannot be specified until the first cross-operator replication run produces data.

\subsection*{Metric 4: Cross-Operator Pathology Replication Rate}

\emph{Definition.} The fraction of named pathologies in \S1.3 independently identified by a naive operator (one who was given the apparatus but not the pathology catalogue) after running the apparatus on a novel corpus.

\emph{What is needed.} A blind replication protocol: operator receives the apparatus (gate library, rubric schema, pre-registration template) but not the pathology catalogue. After running $k$ evaluation cycles on their own corpus, they report which recurring failure modes they observed. An independent rater matches the operator's report against the catalogue.

\emph{Why this metric matters.} The decomposition claim is that the pathologies are ``recurring failure modes of arguments under optimization pressure,'' not artifacts of one system's architecture. A pathology independently rediscovered by a naive operator is confirmed as a real structural pattern. A pathology never independently rediscovered is demoted to ``observed in a single system'' until further evidence arrives. This is the operational form of the falsification condition stated at the end of \S1.4.

\end{document}
